% ============================================================
% Cryptocurrency and Blockchain Technology
% EE465 – IIT Bombay
% Prof. Saravanan Vijayakumaran
% ============================================================
\documentclass[12pt,a4paper,openany]{book}

% ---------- packages ----------
\usepackage{amsmath,amssymb,amsthm}
\usepackage{graphicx}
\usepackage[margin=1in]{geometry}
\usepackage[colorlinks=true,linkcolor=blue,urlcolor=cyan,citecolor=green]{hyperref}
\usepackage{tikz}
\usetikzlibrary{shapes.geometric,arrows.meta,positioning,calc,trees}
\usepackage{listings}
\usepackage{xcolor}
\usepackage{booktabs}
\usepackage{array}
\usepackage{enumitem}
\usepackage[most]{tcolorbox}
\usepackage{fancyhdr}
\usepackage[protrusion=true,expansion=false]{microtype}
\usepackage{mathtools}
\usepackage{mdframed}

% ---------- page style ----------
\pagestyle{fancy}
\fancyhf{}
\fancyhead[LE,RO]{\thepage}
\fancyhead[LO]{\nouppercase{\rightmark}}
\fancyhead[RE]{\nouppercase{\leftmark}}
\renewcommand{\headrulewidth}{0.4pt}

% ---------- theorem environments ----------
\theoremstyle{definition}
\newtheorem{definition}{Definition}[chapter]
\newtheorem{example}{Example}[chapter]
\theoremstyle{plain}
\newtheorem{theorem}{Theorem}[chapter]
\newtheorem{lemma}[theorem]{Lemma}
\newtheorem{corollary}[theorem]{Corollary}
\newtheorem{proposition}[theorem]{Proposition}
\theoremstyle{remark}
\newtheorem*{remark}{Remark}

% ---------- coloured boxes ----------
\tcbuselibrary{skins,breakable}

\newtcolorbox{keyinsight}[1][Key Insight]{
  enhanced, breakable,
  colback=blue!5!white, colframe=blue!60!black,
  fonttitle=\bfseries, title=#1,
  before skip=6pt, after skip=6pt
}

\newtcolorbox{exambox}[1][Exam Note]{
  enhanced, breakable,
  colback=yellow!10!white, colframe=orange!80!black,
  fonttitle=\bfseries, title=#1,
  before skip=6pt, after skip=6pt
}

\newtcolorbox{definitionbox}[1][Definition]{
  enhanced, breakable,
  colback=green!4!white, colframe=green!60!black,
  fonttitle=\bfseries, title=#1,
  before skip=6pt, after skip=6pt
}

% ---------- listings ----------
\definecolor{codebg}{rgb}{0.97,0.97,0.97}
\definecolor{codeframe}{rgb}{0.7,0.7,0.7}
\definecolor{kwcolor}{rgb}{0.0,0.0,0.8}
\definecolor{strcolor}{rgb}{0.6,0.1,0.1}
\definecolor{cmtcolor}{rgb}{0.4,0.4,0.4}

\lstset{
  backgroundcolor=\color{codebg},
  frame=single, rulecolor=\color{codeframe},
  basicstyle=\ttfamily\small,
  keywordstyle=\color{kwcolor}\bfseries,
  stringstyle=\color{strcolor},
  commentstyle=\color{cmtcolor}\itshape,
  numbers=left, numberstyle=\tiny\color{gray},
  numbersep=6pt,
  breaklines=true, breakatwhitespace=true,
  tabsize=2, showstringspaces=false,
  captionpos=b
}

\lstdefinelanguage{Solidity}{
  keywords={pragma,solidity,contract,function,returns,public,private,
            internal,external,view,pure,payable,memory,storage,calldata,
            uint256,uint8,uint,int256,int,bool,address,bytes32,bytes,
            string,mapping,struct,event,emit,require,revert,if,else,
            for,while,return,new,delete,true,false,constructor,modifier,
            interface,library,is,import,using,override,virtual,indexed},
  comment=[l]{//},
  morecomment=[s]{/*}{*/},
  morestring=[b]"
}

\lstdefinelanguage{JavaScript}{
  keywords={let,const,var,function,return,async,await,import,export,
            default,from,class,new,this,typeof,instanceof,if,else,
            for,while,do,break,continue,try,catch,throw,true,false,
            null,undefined,of,in},
  comment=[l]{//},
  morecomment=[s]{/*}{*/},
  morestring=[b]",
  morestring=[b]',
  morestring=[b]`
}

% ---------- misc ----------
\setlength{\parskip}{4pt}

% ============================================================
\begin{document}
% ============================================================

% ---- title page ----
\begin{titlepage}
  \centering
  \vspace*{3cm}
  {\Huge\bfseries Cryptocurrency and\\ Blockchain Technology\par}
  \vspace{0.8cm}
  \rule{\textwidth}{1pt}\par
  \vspace{0.6cm}
  {\LARGE EE\,465 --- IIT Bombay\par}
  \vspace{1.2cm}
  {\large\itshape Prof.\ Saravanan Vijayakumaran\par}
  {\large Department of Electrical Engineering\par}
  {\large Indian Institute of Technology Bombay\par}
  \vspace{1.2cm}
  {\large 2026\par}
  \vspace{2cm}
  \begin{minipage}{0.85\textwidth}
    \itshape
    These lecture notes cover the complete EE\,465 course.  Topics span
    Bitcoin (transactions, mining, scripting), cryptographic foundations
    (hash functions, group theory, elliptic curve cryptography), Ethereum
    (accounts, EVM, smart contracts, Solidity, data structures), privacy
    mechanisms (Tornado Cash, zk-SNARKs), and consensus theory
    (Byzantine broadcast, Tendermint, Ethereum Proof-of-Stake).
  \end{minipage}
\end{titlepage}

\tableofcontents

% ================================================================
\chapter{Introduction to Bitcoin}
\label{chap:bitcoin_intro}
% ================================================================

\section{What Is Bitcoin?}

Bitcoin is a decentralised peer-to-peer electronic cash system introduced
in 2008 by the pseudonymous \textbf{Satoshi Nakamoto}.  The core
innovation is solving the \emph{double-spending problem} without a
trusted central authority.

\begin{keyinsight}
  Bitcoin replaces institutional trust (a bank) with mathematical trust
  (cryptographic proof) enforced by a distributed network of nodes
  running identical validation rules.
\end{keyinsight}

\subsection{The Double-Spending Problem}

In physical cash Alice hands Bob a \$20 note and can no longer spend it.
Digital files can be copied trivially, so Alice could send the same
digital coin to both Bob and Carol.  Traditional systems rely on a bank
to prevent this.  Bitcoin prevents it through the \textbf{blockchain}: a
public ledger whose ordering is agreed upon by \emph{proof of work}.

\subsection{Network Participants}

\begin{description}[leftmargin=2.2cm,style=sameline]
  \item[Full node]   Downloads and validates every block and transaction.
  \item[Miner]       Bundles transactions into blocks and solves PoW.
  \item[Light client] Stores only block headers; uses Merkle proofs.
  \item[Wallet]      Manages keys, creates and broadcasts transactions.
\end{description}

\section{Cryptographic Keys and Addresses}

\subsection{Key Pair}

\begin{itemize}
  \item \textbf{Private key} $k$: uniformly random 256-bit integer in $[1,\,n-1]$.
  \item \textbf{Public key} $K = k\cdot G$: a point on the secp256k1 curve
        (Chapter~\ref{chap:ecc}).
\end{itemize}

\subsection{Bitcoin Address Derivation}

\begin{align}
  H_1 &= \mathrm{SHA256}(K) \\
  H_2 &= \mathrm{RIPEMD160}(H_1) \quad (\text{20 bytes}) \\
  \text{checksum} &= \mathrm{SHA256}(\mathrm{SHA256}(\texttt{0x00}\|H_2))[0:4] \\
  \text{address} &= \mathrm{Base58Check}(\texttt{0x00}\|H_2\|\text{checksum})
\end{align}

\begin{exambox}
  The address is a \emph{hash} of the public key, not the public key
  itself.  The public key is revealed only when spending.  This gives
  limited post-quantum resistance.
\end{exambox}

\section{The Blockchain}

\subsection{Block Header}

\begin{table}[h]
\centering
\caption{Bitcoin Block Header Fields (80 bytes total)}
\begin{tabular}{llp{7cm}}
\toprule
\textbf{Field} & \textbf{Size} & \textbf{Description} \\
\midrule
\texttt{nVersion}      & 4 B  & Protocol version \\
\texttt{hashPrevBlock} & 32 B & SHA256d of previous block header \\
\texttt{hashMerkleRoot}& 32 B & Merkle root of all transactions \\
\texttt{nTime}         & 4 B  & Unix timestamp \\
\texttt{nBits}         & 4 B  & Compact target encoding \\
\texttt{nNonce}        & 4 B  & Variable miners iterate to meet target \\
\bottomrule
\end{tabular}
\end{table}

\begin{figure}[h]
\centering
\begin{tikzpicture}[node distance=1.8cm,
  blk/.style={rectangle,draw,thick,minimum width=2.8cm,
              minimum height=0.9cm,font=\small}]
  \node[blk] (b0) {Block $N\!-\!1$};
  \node[blk,right=of b0] (b1) {Block $N$};
  \node[blk,right=of b1] (b2) {Block $N\!+\!1$};
  \draw[-{Stealth},thick] (b1.west) -- (b0.east)
        node[midway,above,font=\footnotesize]{hashPrevBlock};
  \draw[-{Stealth},thick] (b2.west) -- (b1.east)
        node[midway,above,font=\footnotesize]{hashPrevBlock};
\end{tikzpicture}
\caption{Each block commits to its predecessor via a hash pointer.}
\label{fig:chain}
\end{figure}

\subsection{Merkle Hash Tree}

Transactions are hashed pairwise up a binary tree until a single
\emph{Merkle root} remains:

\begin{align}
  h_{00} &= H(t_0),\quad h_{01} = H(t_1),\quad h_0 = H(h_{00}\|h_{01}) \\
  h_{10} &= H(t_2),\quad h_{11} = H(t_3),\quad h_1 = H(h_{10}\|h_{11}) \\
  \text{root} &= H(h_0\|h_1)
\end{align}

If the transaction count is not a power of two, the last hash is
duplicated to pad the tree.

\textbf{Merkle proofs.}  A light client can verify that $t_i$ is in a
block using $O(\log n)$ sibling hashes, without downloading the full
block.

\subsection{Tamper Resistance and the 51\% Attack}

Modifying any transaction in block $B_N$ changes $H(B_N)$, which breaks
the hash pointer in $B_{N+1}$, forcing a re-mine of every subsequent
block.  An attacker must outpace all honest miners:

\begin{align}
  \text{Current Bitcoin hashrate} &\approx 10^{21}\ \text{H/s} \\
  \text{Cost of 50\% attack} &\approx \$7.5\ \text{billion in hardware}
\end{align}

\begin{keyinsight}
  Bitcoin's security is economic: attacking the network costs more than
  the potential gain, as long as no entity controls $\geq 50\%$ of the
  hashrate.
\end{keyinsight}

\section{The Bitcoin P2P Network}

Bitcoin has no central server.  Every full node maintains a copy of the
entire blockchain ($\approx 600$~GB as of 2024) and communicates
directly with other nodes.

\subsection{Node Discovery}

\begin{enumerate}
  \item A new node starts by connecting to \textbf{DNS seeds} —
        hard-coded domain names maintained by the Bitcoin community that
        return IP addresses of well-connected nodes.
  \item The node sends a \texttt{version} message to its peer, including
        its best block height.
  \item The peer responds with \texttt{verack} and its own
        \texttt{version}.
  \item They exchange \texttt{addr} messages to discover more peers.
  \item A node typically maintains 8 outbound and up to 125 inbound
        connections.
\end{enumerate}

\subsection{Transaction Propagation}

When a user broadcasts a transaction:
\begin{enumerate}
  \item The wallet sends the transaction to a connected full node.
  \item The node validates the transaction (correct signatures, UTXOs
        exist, fee $\geq$ minimum relay fee).
  \item If valid, the node adds it to its \textbf{mempool} and relays
        an \texttt{inv} (inventory) message to its peers.
  \item Peers that do not yet have the transaction respond with
        \texttt{getdata}.
  \item The transaction is forwarded; it reaches most of the network
        within a few seconds.
\end{enumerate}

\subsection{Block Propagation and Orphan Blocks}

When a miner finds a valid block:
\begin{enumerate}
  \item It broadcasts the block immediately.
  \item Other nodes validate (check PoW, check all transactions) and
        relay.
  \item If two miners find valid blocks at the same height almost
        simultaneously, a temporary \textbf{fork} occurs.
  \item Nodes follow the \textbf{longest chain rule} (more precisely,
        the chain with the most cumulative work).
  \item The fork resolves when one branch gets the next block; nodes on
        the shorter branch switch to the longer chain.
  \item Blocks on the abandoned branch are called \textbf{orphan blocks}
        (stale blocks).  Their transactions return to the mempool.
\end{enumerate}

\begin{keyinsight}
  The longest-chain rule is not just about block count — it is about
  cumulative proof-of-work (sum of difficulties).  A single block with
  higher difficulty can outweigh two blocks with lower difficulty.
\end{keyinsight}

\subsection{Simplified Payment Verification (SPV)}

Light clients implement \textbf{SPV}: they download only block headers
(80 bytes each) rather than full blocks.

\begin{itemize}
  \item A block header contains the Merkle root of all transactions.
  \item To verify that transaction $t_i$ is in block $B$, the SPV client
        requests a \textbf{Merkle proof}: the sibling hashes on the path
        from leaf $t_i$ to the root.
  \item The client recomputes the root from $t_i$ and the proof and
        checks it matches the \texttt{hashMerkleRoot} in the header.
  \item This requires $O(\log n)$ hashes instead of downloading all $n$
        transactions.
  \item SPV trusts that the chain with the most PoW contains valid
        transactions — it does not independently verify all transactions.
\end{itemize}

\section{The Genesis Block}

The \textbf{genesis block} (block 0) was mined by Satoshi Nakamoto on
January 3, 2009.  It is hard-coded into the Bitcoin software.

\begin{itemize}
  \item The coinbase transaction contains the famous message:
        \emph{``The Times 03/Jan/2009 Chancellor on brink of second
        bailout for banks.''} — a newspaper headline proving it was not
        pre-mined before that date.
  \item Its \texttt{hashPrevBlock} is all zeros (no predecessor).
  \item The 50~BTC reward in the genesis block is unspendable due to an
        original quirk in the Bitcoin software.
\end{itemize}

\section{The UTXO Set}

At any point in time, the \textbf{UTXO set} is the complete collection
of all unspent transaction outputs.

\begin{itemize}
  \item Every full node maintains the UTXO set in memory (or on a fast
        database like LevelDB).
  \item As of 2024, the UTXO set contains $\approx 100$ million UTXOs,
        requiring $\approx 5$~GB of RAM.
  \item When a transaction is processed:
        \begin{enumerate}
          \item Each input removes a UTXO from the set.
          \item Each output adds a new UTXO to the set.
        \end{enumerate}
  \item A transaction is valid only if all its inputs reference UTXOs
        currently in the UTXO set (and are properly signed).
\end{itemize}

\begin{exambox}
  \textbf{Why not an account balance?}  The UTXO model enables parallel
  transaction validation (non-conflicting UTXOs can be validated
  simultaneously) and makes the state explicit: you can prove you have
  coins by pointing to specific UTXOs, without trusting a running total.
\end{exambox}

\subsection{Key Takeaways}

\begin{itemize}
  \item Prevents double spending and tampering via blockchain + PoW.
  \item Secure only if no party controls $\geq 50\%$ of hashrate.
  \item Mining difficulty adjusts every 2016 blocks to target 10~min/block.
  \item Block subsidy halves every 210\,000 blocks; total supply capped at
        21\,000\,000~BTC.
  \item The P2P network propagates transactions and blocks without any
        central server.
  \item The UTXO model is the fundamental accounting primitive: no
        account balances, only unspent outputs.
\end{itemize}


% ================================================================
\chapter{Cryptographic Hash Functions}
\label{chap:hash}
% ================================================================

\section{Definition}

\begin{definitionbox}
  A \textbf{cryptographic hash function} is a map
  $H:\{0,1\}^*\!\to\{0,1\}^n$ satisfying:
  \begin{enumerate}
    \item \textbf{Efficiency}: $H(x)$ is fast to compute for any $x$.
    \item \textbf{Preimage resistance}: Given $h$, it is infeasible to
          find any $x$ with $H(x)=h$.
    \item \textbf{Second-preimage resistance}: Given $x$, it is infeasible
          to find $x'\neq x$ with $H(x')=H(x)$.
    \item \textbf{Collision resistance}: It is infeasible to find any
          distinct pair $x\neq x'$ with $H(x)=H(x')$.
  \end{enumerate}
\end{definitionbox}

\subsection{Security Hierarchy}

Collision resistance $\Rightarrow$ Second-preimage resistance (roughly).
The reverse does not hold in general.

\subsection{Attack Complexities}

\begin{table}[h]
\centering
\caption{Attack complexities for an $n$-bit hash}
\begin{tabular}{lll}
\toprule
\textbf{Attack} & \textbf{Complexity} & \textbf{SHA-256 ($n=256$)} \\
\midrule
Brute-force preimage & $O(2^n)$ & $2^{256}$ \\
Birthday collision   & $O(2^{n/2})$ & $2^{128}$ \\
\bottomrule
\end{tabular}
\end{table}

\textbf{Birthday paradox}: In a group of $\approx 1.2\sqrt{N}$ people,
there is a $>50\%$ chance two share a birthday.  For hash collisions,
$\approx 2^{n/2}$ evaluations suffice.

\section{SHA-256}

SHA-256 (Secure Hash Algorithm, 256-bit output) is a
\textbf{Merkle--Damg\aa rd} construction:
\begin{enumerate}
  \item Pad message to a multiple of 512 bits.
  \item Split into 512-bit blocks $M_1,\ldots,M_k$.
  \item Iterate: $H_i = \text{compress}(H_{i-1},M_i)$ with $H_0$ = fixed IV.
  \item Output $H_k$ (256 bits).
\end{enumerate}

The compression function performs 64 rounds of mixing using:
bitwise operations (AND, OR, XOR, NOT), bit rotations/shifts, modular
addition ($\bmod 2^{32}$), and a schedule of 64 derived 32-bit words.

\subsection{Bitcoin Usage of SHA-256}

\begin{table}[h]
\centering
\caption{SHA-256 in Bitcoin}
\begin{tabular}{ll}
\toprule
\textbf{Purpose} & \textbf{Operation} \\
\midrule
Block hash / PoW & $\mathrm{SHA256}(\mathrm{SHA256}(\text{header}))$ \\
Transaction ID & $\mathrm{SHA256}(\mathrm{SHA256}(\text{serialised tx}))$ \\
Merkle tree nodes & $\mathrm{SHA256}(\mathrm{SHA256}(L\|R))$ \\
Address step 1 & $\mathrm{SHA256}(\text{pubkey})$ \\
Checksum & First 4 B of $\mathrm{SHA256}(\mathrm{SHA256}(\text{payload}))$ \\
\bottomrule
\end{tabular}
\end{table}

\section{SHA-256 Internal Structure in Detail}

\subsection{Padding}

Before hashing, the message $M$ of bit-length $L$ is padded to a
multiple of 512 bits:
\begin{enumerate}
  \item Append a single bit `1'.
  \item Append `0' bits until the length $\equiv 448 \pmod{512}$.
  \item Append $L$ as a 64-bit big-endian integer.
\end{enumerate}

\textbf{Example}: hashing the 3-byte string \texttt{"abc"} ($L=24$~bits):
\begin{lstlisting}[language={}]
Message:   61 62 63
Padding:   80 00...00 00 00 00 00 00 00 00 18
                  (length = 24 = 0x18)
Total: 512 bits (one block)
\end{lstlisting}

\subsection{The Compression Function}

SHA-256 uses eight 32-bit state words $H_0,\ldots,H_7$ initialised to
the fractional parts of the square roots of the first 8 primes:
\[
H_0 = \texttt{6a09e667},\quad H_1 = \texttt{bb67ae85},\quad\ldots
\]

For each 512-bit message block, 64 rounds are performed.  Each round
uses one of 64 constants $K_i$ (derived from cube roots of the first 64
primes) and the \textbf{message schedule} $W_0,\ldots,W_{63}$.

\textbf{Message schedule} ($W_i$ for $i\geq 16$):
\[
W_i = \sigma_1(W_{i-2}) + W_{i-7} + \sigma_0(W_{i-15}) + W_{i-16}
\]
where $\sigma_0(x)=\mathrm{ROTR}^7(x)\oplus\mathrm{ROTR}^{18}(x)\oplus\mathrm{SHR}^3(x)$
and $\sigma_1$ is similar with rotation amounts 17, 19, 10.

\textbf{Round function}: let $a,b,c,d,e,f,g,h$ be the working variables
(initialised from $H_0,\ldots,H_7$):
\begin{align}
  T_1 &= h + \Sigma_1(e) + \mathrm{Ch}(e,f,g) + K_i + W_i \\
  T_2 &= \Sigma_0(a) + \mathrm{Maj}(a,b,c) \\
  h&\gets g,\; g\gets f,\; f\gets e,\; e\gets d+T_1 \\
  d&\gets c,\; c\gets b,\; b\gets a,\; a\gets T_1+T_2
\end{align}

\begin{align*}
  \mathrm{Ch}(x,y,z)  &= (x\wedge y)\oplus(\neg x\wedge z) \\
  \mathrm{Maj}(x,y,z) &= (x\wedge y)\oplus(x\wedge z)\oplus(y\wedge z) \\
  \Sigma_0(x)         &= \mathrm{ROTR}^{2}(x)\oplus\mathrm{ROTR}^{13}(x)\oplus\mathrm{ROTR}^{22}(x) \\
  \Sigma_1(x)         &= \mathrm{ROTR}^{6}(x)\oplus\mathrm{ROTR}^{11}(x)\oplus\mathrm{ROTR}^{25}(x)
\end{align*}

After 64 rounds, the block output is added back to $H_0,\ldots,H_7$
modulo $2^{32}$.

\subsection{Avalanche Effect}

A one-bit change in the input propagates through the round function to
affect (roughly) half the output bits within a few rounds.

\begin{example}
\begin{lstlisting}[language={}]
SHA256("abc")    = ba7816bf 8f01cfea 414140de 5dae2ec7
                   3b338c12 12133393 8e04d2eb 4bb6e54
SHA256("abd")    = 3909e2e0 5f6dfb8b 6a6e1a3a 7b1a0b69
                   e9e4f5c3 9b0e7aa4 f2c2e75b e8f5a9c2
\end{lstlisting}
One character changes (c$\to$d); every output byte is completely different.
\end{example}

\subsection{Length Extension Attacks}

SHA-256's Merkle-Damg\aa rd construction is vulnerable to \textbf{length
extension attacks}: given $H(m)$, an attacker can compute
$H(m\|\text{padding}\|m')$ without knowing $m$.

This is why Bitcoin uses \textbf{SHA256d} = SHA-256(SHA-256($\cdot$)):
applying SHA-256 twice breaks the Merkle-Damg\aa rd chain and eliminates
length extension vulnerability.

\section{RIPEMD-160}

RIPEMD-160 produces 20-byte output.  Combined with SHA-256, it
compresses a public key hash to 20 bytes for use in Bitcoin addresses.
Using two independent hash families gives defence-in-depth against
weaknesses in either.

\section{Hash Pointers}

A \textbf{hash pointer} stores both a data pointer and $H(\text{data})$.
Modifying the data invalidates the pointer.  Bitcoin's blockchain is a
linked list of hash pointers:

\begin{align}
  \boxed{\texttt{hashPrevBlock} = \mathrm{SHA256d}(\text{previous header})}
\end{align}


% ================================================================
\chapter{Group Theory and Modular Arithmetic}
\label{chap:groups}
% ================================================================

\section{Groups}

\begin{definitionbox}
  A \textbf{group} $(G,\cdot)$ is a set $G$ with a binary operation
  satisfying:
  \begin{enumerate}[label=(\roman*)]
    \item \textbf{Closure}: $a\cdot b\in G$ for all $a,b\in G$.
    \item \textbf{Associativity}: $(a\cdot b)\cdot c = a\cdot(b\cdot c)$.
    \item \textbf{Identity}: $\exists\,e$ s.t.\ $e\cdot a=a\cdot e=a$ for all $a$.
    \item \textbf{Inverses}: $\forall\,a\;\exists\,a^{-1}$ s.t.\
          $a\cdot a^{-1}=e$.
  \end{enumerate}
  If also $a\cdot b=b\cdot a$ for all $a,b$, the group is \textbf{abelian}.
\end{definitionbox}

\subsection{Examples Relevant to Cryptography}

\begin{table}[h]
\centering
\caption{Common cryptographic groups}
\begin{tabular}{llll}
\toprule
\textbf{Group} & \textbf{Operation} & \textbf{Order} & \textbf{Cyclic?}\\
\midrule
$(\mathbb{Z}_n,+)$ & Addition mod $n$ & $n$ & Yes \\
$(\mathbb{Z}_p^*,\times)$ & Multiplication mod prime $p$ & $p-1$ & Yes \\
$E(\mathbb{F}_p)$ & Point addition on curve & $\approx p$ & Yes (usually) \\
\bottomrule
\end{tabular}
\end{table}

\section{Cyclic Groups and the Discrete Logarithm}

\begin{definitionbox}
  $G$ is \textbf{cyclic} if $\exists\,g\in G$ (a \emph{generator}) such that
  $G=\{g^0,g^1,\ldots,g^{|G|-1}\}$.
\end{definitionbox}

\textbf{Exponentiation} (easy): given $k$, compute $h=g^k$ in
$O(\log k)$ group operations via double-and-add.

\textbf{Discrete Logarithm Problem (DLP)}: given $h=g^k$, find $k$.
Believed hard for suitable groups.

\section{Modular Arithmetic}

\[
a \equiv b\pmod n \iff n\mid(a-b)
\]

\textbf{Modular inverse}: if $\gcd(a,n)=1$ then $a^{-1}\bmod n$ exists,
computable via the Extended Euclidean Algorithm.

\begin{theorem}[Fermat's Little Theorem]
  For prime $p$ and $\gcd(a,p)=1$:
  \[
    \boxed{a^{p-1}\equiv 1\pmod p} \implies a^{-1}\equiv a^{p-2}\pmod p
  \]
\end{theorem}

\section{Finite Fields}

\begin{definitionbox}
  $\mathbb{F}_p=\mathbb{Z}_p$ for prime $p$: the set $\{0,1,\ldots,p-1\}$
  with addition and multiplication modulo $p$.  Every non-zero element has
  a multiplicative inverse.
\end{definitionbox}

Elliptic curves in Bitcoin are defined over $\mathbb{F}_p$ for a 256-bit prime $p$.

\subsection{The Extended Euclidean Algorithm}

To find $a^{-1}\bmod n$, use the Extended Euclidean Algorithm (EEA):

\begin{example}
  Find $7^{-1}\bmod 11$:

  Run $\gcd(7,11)$:
  \[
    11 = 1\cdot 7 + 4,\quad 7 = 1\cdot 4+3,\quad 4=1\cdot 3+1,\quad
    3=3\cdot 1
  \]
  Back-substitute:
  \[
    1 = 4 - 1\cdot 3 = 4 - 1\cdot(7-4) = 2\cdot 4 - 7
      = 2\cdot(11-7) - 7 = 2\cdot 11 - 3\cdot 7
  \]
  So $-3\cdot 7\equiv 1\pmod{11}$, i.e.\ $7^{-1}\equiv -3\equiv 8\pmod{11}$.

  Verify: $7\times 8 = 56 = 5\times 11 + 1\equiv 1\pmod{11}$. \checkmark
\end{example}

\subsection{Worked Example: Discrete Logarithm in $\mathbb{Z}_{11}^*$}

$\mathbb{Z}_{11}^* = \{1,2,3,4,5,6,7,8,9,10\}$ is a cyclic group of
order 10.

\textbf{Finding generators}: $g$ is a generator iff $\mathrm{ord}(g)=10$.

\begin{table}[h]
\centering
\caption{Powers of $g=2$ in $\mathbb{Z}_{11}^*$}
\begin{tabular}{cccccccccc}
\toprule
$2^1$ & $2^2$ & $2^3$ & $2^4$ & $2^5$ & $2^6$ & $2^7$ & $2^8$ & $2^9$ & $2^{10}$ \\
\midrule
2 & 4 & 8 & 5 & 10 & 9 & 7 & 3 & 6 & 1 \\
\bottomrule
\end{tabular}
\end{table}

$g=2$ generates all of $\mathbb{Z}_{11}^*$, so it is a generator.

\textbf{DLP example}: Given $h=7$, find $k$ such that $2^k\equiv 7\pmod{11}$.
From the table: $2^7=7$, so $k=7$.

In practice (with 256-bit groups), no efficient algorithm is known.

\section{Lagrange's Theorem}

\begin{theorem}
  For finite group $G$ and subgroup $H\leq G$, $|H|$ divides $|G|$.
\end{theorem}

\textbf{Corollary}: The order of any $g\in G$ divides $|G|$, so
$g^{|G|}=e$.

\section{Quadratic Residues}

$a\in\mathbb{F}_p^*$ is a \textbf{quadratic residue (QR)} if
$\exists\,x$ with $x^2\equiv a$.

\begin{theorem}[Euler's Criterion]
  \[
    a^{(p-1)/2}\equiv\begin{cases}1\pmod p & a \text{ is a QR}\\
    -1\pmod p & a \text{ is a QNR}\end{cases}
  \]
\end{theorem}

Exactly $(p-1)/2$ elements of $\mathbb{F}_p^*$ are QRs.


% ================================================================
\chapter{Elliptic Curve Cryptography in Bitcoin}
\label{chap:ecc}
% ================================================================

\section{Elliptic Curves}

\begin{definitionbox}
  An \textbf{elliptic curve} over $\mathbb{F}_p$ is the set of points
  $(x,y)\in\mathbb{F}_p^2$ satisfying
  \begin{equation}
    y^2 \equiv x^3+ax+b \pmod p, \quad \Delta=-16(4a^3+27b^2)\neq 0,
  \end{equation}
  together with a formal \emph{point at infinity} $\mathcal{O}$.
\end{definitionbox}

\subsection{Bitcoin's Curve: secp256k1}

$a=0$, $b=7$:
\begin{equation}
  \boxed{y^2 \equiv x^3+7 \pmod p}
\end{equation}
\[
  p = 2^{256}-2^{32}-977
  = \texttt{FFFFFFFFFFFFFFFFFFFFFFFFFFFFFFFFFFFFFFFFFFFFFFFFFFFFEFFFFFC2F}_{16}
\]

Group order:
\[
  n = \texttt{FFFFFFFFFFFFFFFFFFFFFFFFFFFFFFFEBAAEDCE6AF48A03BBFD25E8CD0364141}_{16}
\]

\section{Group Law on Elliptic Curves}

The set $E(\mathbb{F}_p)$ forms an abelian group under \emph{point addition}.

\subsection{Point Addition: $P+Q$ with $P\neq\pm Q$}

\begin{align}
  \lambda &= \frac{y_2-y_1}{x_2-x_1}\pmod p \\
  x_3 &= \lambda^2-x_1-x_2 \pmod p \\
  y_3 &= \lambda(x_1-x_3)-y_1 \pmod p
\end{align}

Geometrically: draw the line through $P$ and $Q$, reflect the third
intersection point over the $x$-axis.

\subsection{Point Doubling: $P+P=2P$}

\begin{align}
  \lambda &= \frac{3x_1^2+a}{2y_1}\pmod p \quad (a=0\text{ for secp256k1})\\
  x_3 &= \lambda^2-2x_1 \pmod p \\
  y_3 &= \lambda(x_1-x_3)-y_1 \pmod p
\end{align}

\subsection{Identity and Inverse}

\begin{itemize}
  \item $P+\mathcal{O}=P$ \quad (identity is $\mathcal{O}$)
  \item $-P=(x,-y\bmod p)$; \quad $P+(-P)=\mathcal{O}$
\end{itemize}

\section{Scalar Multiplication}

\[
  kP = \underbrace{P+P+\cdots+P}_{k}
\]

Efficient via \textbf{double-and-add} in $O(\log k)$ operations:

\begin{lstlisting}[language={}]
R = O (point at infinity)
for bit b in binary(k) from MSB to LSB:
    R = 2R         (doubling)
    if b == 1:
        R = R + P  (addition)
return R
\end{lstlisting}

\section{Worked Example: Point Addition over $\mathbb{F}_{17}$}

Consider the curve $E: y^2\equiv x^3+2x+2\pmod{17}$.

\begin{example}
  Let $P=(5,1)$ and $Q=(6,3)$ (both on the curve; verify:
  $1^2=1\equiv 5^3+2\cdot5+2=137\equiv 1\pmod{17}$~\checkmark).

  \textbf{Compute $P+Q$}:
  \begin{align*}
    \lambda &= \frac{y_2-y_1}{x_2-x_1} = \frac{3-1}{6-5} = \frac{2}{1} = 2 \\
    x_3 &= 2^2 - 5 - 6 = 4-11 = -7\equiv 10\pmod{17} \\
    y_3 &= 2(5-10)-1 = -10-1 = -11\equiv 6\pmod{17}
  \end{align*}
  So $P+Q = (10,6)$.
\end{example}

\begin{example}
  \textbf{Compute $2P$} with $P=(5,1)$:
  \begin{align*}
    \lambda &= \frac{3\cdot5^2+2}{2\cdot1} = \frac{77}{2}
             = \frac{9}{2}\pmod{17}
             = 9\cdot 2^{-1}\pmod{17} \\
    2^{-1}&\equiv 9\pmod{17} \quad(\text{since }2\cdot 9=18\equiv 1) \\
    \lambda &= 9\cdot 9 = 81\equiv 13\pmod{17} \\
    x_3 &= 13^2 - 2\cdot5 = 169-10 = 159\equiv 6\pmod{17} \\
    y_3 &= 13(5-6)-1 = -13-1=-14\equiv 3\pmod{17}
  \end{align*}
  So $2P=(6,3)=Q$.  Hence $P+Q = 2P+P=3P$, meaning $Q=2P$.
\end{example}

\begin{example}
  \textbf{Scalar multiplication $4P$}:
  \[
    2P=(6,3),\quad 4P=2(2P)=2(6,3)
  \]
  Compute $2(6,3)$: $\lambda=(3\cdot36+2)/(2\cdot3)=110/6$.
  $6^{-1}\equiv 3\pmod{17}$ (since $6\cdot3=18\equiv1$).
  $\lambda=110\cdot3\bmod17=330\bmod17=6$.
  $x_3=36-12=24\equiv7$, $y_3=6(6-7)-3=-6-3=-9\equiv8\pmod{17}$.
  So $4P=(7,8)$.
\end{example}

The point at infinity $\mathcal{O}$ is the additive identity.
The curve over $\mathbb{F}_{17}$ has 18 points (including $\mathcal{O}$),
forming a cyclic group of order 18.

\section{Elliptic Curve Discrete Logarithm Problem (ECDLP)}

\begin{definitionbox}
  Given generator $G$ and point $Q=kG$, the \textbf{ECDLP} is to find $k$.
\end{definitionbox}

Best known attacks (Pollard's $\rho$) require $O(\sqrt{n})$ operations.
For secp256k1: $\sqrt{n}\approx 2^{128}$ — computationally infeasible.

\section{Key Generation}

\begin{align}
  k &\xleftarrow{\$} \{1,\ldots,n-1\} \quad\text{(private key)} \\
  \boxed{K} &= \boxed{k\cdot G} \quad\text{(public key, 33 bytes compressed)}
\end{align}

\textbf{Compressed public key}: prefix \texttt{02} if $y$ is even,
\texttt{03} if $y$ is odd, followed by $x$-coordinate (32 bytes).

\section{ECDSA: Elliptic Curve Digital Signature Algorithm}

\subsection{Signing}

Given message $m$, private key $k$, generator $G$, group order $n$:
\begin{enumerate}
  \item $e = H(m)$ \quad (hash of message)
  \item Choose random nonce $r\in[1,n-1]$
  \item $R = r\cdot G = (x_R,y_R)$; set $r_{\text{sig}} = x_R\bmod n$
  \item $s = r^{-1}(e + k\cdot r_{\text{sig}})\bmod n$
  \item Signature: $(r_{\text{sig}},s)$
\end{enumerate}

\begin{equation}
  \boxed{s = r^{-1}\!\bigl(H(m)+k\cdot r_{\text{sig}}\bigr)\bmod n}
\end{equation}

\subsection{Verification}

Given $(r_{\text{sig}},s)$, public key $K$:
\begin{enumerate}
  \item $e = H(m)$
  \item $u_1 = s^{-1}e\bmod n$, \quad $u_2 = s^{-1}r_{\text{sig}}\bmod n$
  \item $R' = u_1 G + u_2 K$
  \item Accept iff $R'_x\equiv r_{\text{sig}}\pmod n$
\end{enumerate}

\begin{exambox}[Security Warning]
  \textbf{Never reuse the nonce $r$!}  If $r$ is reused for two messages
  $m_1,m_2$, the private key can be recovered:
  $k = (s_1-s_2)^{-1}(H(m_1)-H(m_2))\bmod n$.
  This caused the PlayStation~3 key leak (2010).
\end{exambox}

\section{ECDSA Worked Conceptual Example}

\begin{example}
  \textbf{Why the verification equation works.}

  We have signature $(r_{\text{sig}}, s)$ where
  $s = r^{-1}(e + k\cdot r_{\text{sig}})\bmod n$.

  During verification we compute $R' = u_1 G + u_2 K$ where
  $u_1 = s^{-1}e\bmod n$ and $u_2 = s^{-1}r_{\text{sig}}\bmod n$.

  \begin{align*}
    R' &= u_1 G + u_2 K = s^{-1}e\cdot G + s^{-1}r_{\text{sig}}\cdot kG \\
       &= s^{-1}(e + k\cdot r_{\text{sig}})\cdot G = s^{-1}\cdot s r\cdot G = r\cdot G = R
  \end{align*}

  So $R'_x = R_x = r_{\text{sig}}$.  The equation works because only
  someone knowing $k$ can produce an $s$ that satisfies it.
\end{example}

\subsection{Schnorr Signatures (Taproot)}

Bitcoin's Taproot upgrade (BIP\,340, activated November 2021) added
\textbf{Schnorr signatures} as an alternative to ECDSA:

\textbf{Signing}: Choose random nonce $r$, compute $R=r\cdot G$.
\[
  s = r + H(R\|K\|m)\cdot k \pmod n
\]
Signature: $(R, s)$.

\textbf{Verification}: Check $s\cdot G = R + H(R\|K\|m)\cdot K$.

\textbf{Advantages over ECDSA}:
\begin{itemize}
  \item \textbf{Linearity}: $n$ Schnorr signatures on the same message
        can be aggregated into a single signature (MuSig).  This enables
        efficient $n$-of-$n$ multisig and reduces blockchain footprint.
  \item \textbf{Simpler security proof}: Provably secure in the random
        oracle model via a tighter reduction.
  \item \textbf{Batch verification}: $n$ signatures can be verified
        faster than $n$ independent verifications.
\end{itemize}

\section{Parameter Summary}

\begin{table}[h]
\centering
\caption{secp256k1 at a glance}
\begin{tabular}{ll}
\toprule
Curve & $y^2=x^3+7\pmod p$ \\
Field prime $p$ & $2^{256}-2^{32}-977$ (256 bits) \\
Group order $n$ & $\approx 2^{256}$ \\
Security level & $\approx 128$ bits \\
Private key & 32 bytes (256 bits) \\
Compressed public key & 33 bytes \\
Uncompressed public key & 65 bytes \\
Signature (DER) & 71--72 bytes \\
\bottomrule
\end{tabular}
\end{table}


% ================================================================
\chapter{Bitcoin Transactions}
\label{chap:btctx}
% ================================================================

\section{The UTXO Model}

\begin{definitionbox}
  A \textbf{UTXO} (Unspent Transaction Output) is an output of a
  previous transaction that has not yet been spent.  A wallet's
  \emph{balance} is the sum of all UTXOs locked to its addresses.
\end{definitionbox}

Unlike the account model (Ethereum), Bitcoin has no notion of a running
balance.  Every transaction explicitly consumes UTXOs as inputs and
creates new UTXOs as outputs.

\begin{keyinsight}
  A UTXO is like a physical coin: you spend it whole.  If you pay 0.5~BTC
  from a 1~BTC UTXO, you create a 0.5~BTC output to the recipient and a
  change output back to yourself.  The difference is the miner fee.
\end{keyinsight}

\section{Transaction Structure}

\begin{table}[h]
\centering
\caption{Bitcoin Transaction Fields}
\begin{tabular}{lll}
\toprule
\textbf{Field} & \textbf{Size} & \textbf{Description} \\
\midrule
version     & 4 B & Format version \\
vin\_count  & varint & Number of inputs \\
vin[]       & variable & Array of inputs \\
vout\_count & varint & Number of outputs \\
vout[]      & variable & Array of outputs \\
locktime    & 4 B & Earliest block/time for inclusion \\
\bottomrule
\end{tabular}
\end{table}

\subsection{Input Fields}

\begin{table}[h]
\centering
\caption{Transaction Input Fields}
\begin{tabular}{lll}
\toprule
\textbf{Field} & \textbf{Size} & \textbf{Description} \\
\midrule
txid       & 32 B & TXID of the transaction being spent \\
vout       & 4 B & Index of the output within that tx \\
scriptSig  & variable & Unlocking script \\
sequence   & 4 B & RBF / relative locktime \\
\bottomrule
\end{tabular}
\end{table}

\subsection{Output Fields}

\begin{table}[h]
\centering
\caption{Transaction Output Fields}
\begin{tabular}{lll}
\toprule
\textbf{Field} & \textbf{Size} & \textbf{Description} \\
\midrule
value       & 8 B & Amount in satoshis ($10^8\ \text{sat}=1\ \text{BTC}$) \\
scriptPubKey & variable & Locking script \\
\bottomrule
\end{tabular}
\end{table}

\textbf{Implicit fee}:
\[
  \text{fee} = \sum\text{input values} - \sum\text{output values}
\]

\section{Bitcoin Script}

Bitcoin uses a \textbf{stack-based} scripting language.

\begin{itemize}
  \item \textbf{scriptPubKey} (locking script): conditions to spend the output.
  \item \textbf{scriptSig} (unlocking script): data that satisfies conditions.
\end{itemize}

A UTXO is spendable iff \texttt{scriptSig}\,$\|$\,\texttt{scriptPubKey}
leaves \texttt{TRUE} on the stack.

\subsection{Key Opcodes}

\begin{table}[h]
\centering
\caption{Common Bitcoin Script opcodes}
\begin{tabular}{ll}
\toprule
\textbf{Opcode} & \textbf{Effect} \\
\midrule
\texttt{OP\_DUP}          & Duplicate top stack item \\
\texttt{OP\_HASH160}      & Apply SHA256 then RIPEMD160 \\
\texttt{OP\_EQUAL}        & Pop two items; push 1 if equal \\
\texttt{OP\_EQUALVERIFY}  & Like EQUAL but fail if not equal \\
\texttt{OP\_CHECKSIG}     & Verify ECDSA signature \\
\texttt{OP\_CHECKMULTISIG}& Verify $m$-of-$n$ multisignature \\
\texttt{OP\_RETURN}       & Mark output unspendable (data) \\
\bottomrule
\end{tabular}
\end{table}

\section{P2PKH — Pay to Public Key Hash}

Most common transaction type; locks output to a Bitcoin address.

\textbf{scriptPubKey}:
\begin{lstlisting}[language={}]
OP_DUP OP_HASH160 <PubKeyHash> OP_EQUALVERIFY OP_CHECKSIG
\end{lstlisting}

\textbf{scriptSig}:
\begin{lstlisting}[language={}]
<Signature> <PublicKey>
\end{lstlisting}

\textbf{Execution trace}:
\begin{enumerate}
  \item Push \texttt{<Sig>}, \texttt{<PubKey>}.
  \item \texttt{OP\_DUP}: duplicate \texttt{<PubKey>}.
  \item \texttt{OP\_HASH160}: hash the copy $\to$ \texttt{<PubKeyHash'>}.
  \item Push \texttt{<PubKeyHash>} from scriptPubKey.
  \item \texttt{OP\_EQUALVERIFY}: assert \texttt{Hash'} $=$ \texttt{Hash}.
  \item \texttt{OP\_CHECKSIG}: verify ECDSA signature. Stack: \texttt{[TRUE]}.
\end{enumerate}

\section{P2SH — Pay to Script Hash}

Allows complex conditions (multisig, timelocks) to be encoded in a
\emph{redeemScript}; the locking script only contains the hash.

\textbf{scriptPubKey}:
\begin{lstlisting}[language={}]
OP_HASH160 <RedeemScriptHash> OP_EQUAL
\end{lstlisting}

\textbf{scriptSig}:
\begin{lstlisting}[language={}]
<Data...> <RedeemScript>
\end{lstlisting}

\section{Multisignature ($m$-of-$n$)}

\begin{lstlisting}[language={}]
// 2-of-3 redeemScript:
OP_2 <PubKey1> <PubKey2> <PubKey3> OP_3 OP_CHECKMULTISIG

// Unlocking:
OP_0 <Sig1> <Sig2>
\end{lstlisting}

\texttt{OP\_0} compensates for a historical off-by-one bug in
\texttt{OP\_CHECKMULTISIG}.

\section{Worked Example: A Complete Bitcoin Transaction}

\begin{example}
  \textbf{Scenario}: Alice wants to pay Bob 0.3~BTC.  Alice has two
  UTXOs:
  \begin{itemize}
    \item UTXO~A: 0.4~BTC from transaction \texttt{txid\,=\,aaa...}, output 0
    \item UTXO~B: 0.2~BTC from transaction \texttt{txid\,=\,bbb...}, output 1
  \end{itemize}
  She decides to spend UTXO~A only (0.4~BTC $>$ 0.3~BTC payment).

  \textbf{Transaction structure}:
  \begin{itemize}
    \item \textbf{Input}: UTXO~A (\texttt{aaa...}:0) with Alice's
          signature and public key in \texttt{scriptSig}.
    \item \textbf{Output 0}: 0.3~BTC locked to Bob's address
          (P2PKH: \texttt{OP\_DUP OP\_HASH160 <Bob\_hash> OP\_EQUALVERIFY OP\_CHECKSIG}).
    \item \textbf{Output 1}: 0.095~BTC locked to Alice's change address.
    \item \textbf{Implicit fee}: $0.4 - 0.3 - 0.095 = 0.005$~BTC paid to miner.
  \end{itemize}

  \textbf{What happens to Alice's UTXO set?}
  \begin{itemize}
    \item UTXO~A (0.4~BTC) is \emph{removed} from the set (spent).
    \item Two new UTXOs are \emph{added}: Bob's 0.3~BTC and Alice's change 0.095~BTC.
    \item UTXO~B (0.2~BTC) is unaffected.
  \end{itemize}
\end{example}

\section{P2WPKH: Pay-to-Witness-Public-Key-Hash}

The SegWit equivalent of P2PKH.  The \texttt{scriptPubKey} is:
\begin{lstlisting}[language={}]
OP_0 <20-byte-pubkey-hash>
\end{lstlisting}

The witness data (signature + public key) is stored in the
\texttt{witness} field, outside the transaction body used to compute the
TXID.  This makes the TXID immutable with respect to signature changes
(fixing transaction malleability).

\section{Timelock Transactions}

\textbf{Absolute timelocks}: \texttt{locktime} field specifies the
earliest block height or timestamp when the transaction can be mined.

\textbf{Relative timelocks} (BIP\,68): the \texttt{sequence} field
specifies a minimum number of blocks or time since the referenced UTXO
was created.

\textbf{OP\_CHECKLOCKTIMEVERIFY (OP\_CLTV)}: a script opcode that fails
if the current block height/time is less than the value on the stack.

\textbf{Use case}: Hash Time-Locked Contracts (HTLCs) — the basis of
the Lightning Network payment channels.  Alice pays Bob only if Bob
reveals a secret $s$ such that $H(s)=h$ before a deadline.

\begin{lstlisting}[language={}]
// HTLC locking script:
OP_IF
    OP_SHA256 <h> OP_EQUALVERIFY OP_DUP OP_HASH160 <Bob_hash>
OP_ELSE
    <deadline> OP_CHECKLOCKTIMEVERIFY OP_DROP OP_DUP OP_HASH160 <Alice_hash>
OP_ENDIF
OP_EQUALVERIFY OP_CHECKSIG
\end{lstlisting}

\section{SegWit and Transaction Malleability}

SegWit (BIP\,141, activated August 2017) moves witness data (signatures)
outside the transaction ID hash:

\begin{itemize}
  \item Fixes \textbf{transaction malleability} (TXID could be changed by
        modifying signatures).
  \item Witness data discounted 75\% for fee calculation.
  \item Enables the \textbf{Lightning Network}.
\end{itemize}

\section{Coinbase Transaction}

Every block's first transaction is the \textbf{coinbase}:
\begin{itemize}
  \item Has no real inputs (creates new coins ex nihilo).
  \item Pays miner: block subsidy + sum of all transaction fees.
  \item Can embed arbitrary data in a special coinbase field.
\end{itemize}

\textbf{Block subsidy schedule} (halving every 210\,000 blocks):
\[
  \text{Total supply} = \sum_{k=0}^{\infty}210000\times\frac{50}{2^k}
  = 21{,}000{,}000\ \text{BTC}
\]


% ================================================================
\chapter{Mining Miscellanea}
\label{chap:mining}
% ================================================================

\section{Proof of Work}

\begin{equation}
  \boxed{\mathrm{SHA256d}(\text{block header}) < T}
\end{equation}

Miners increment \texttt{nNonce} (and extra nonce in the coinbase) until
the block hash falls below target $T$.

\subsection{Difficulty and Adjustment}

\[
  D = \frac{T_{\text{genesis}}}{T_{\text{current}}}
\]

Every 2016 blocks ($\approx 2$ weeks) the target adjusts:
\[
  T_{\text{new}} = T_{\text{old}}
  \times\frac{\text{actual time for 2016 blocks}}{2016\times10\,\text{min}}
  \quad\text{(capped at factor of 4)}
\]

\subsection{Expected Hashes per Block}

If each hash independently satisfies the target with probability
$p=T/2^{256}$, the expected number of hashes is $1/p$.

\section{Mining Hardware}

\begin{table}[h]
\centering
\caption{Mining hardware evolution}
\begin{tabular}{llll}
\toprule
\textbf{Era} & \textbf{Hardware} & \textbf{Typical rate} & \textbf{Efficiency}\\
\midrule
2009--10 & CPU  & 10~MH/s     & Poor \\
2010--11 & GPU  & 100~MH/s    & Better \\
2011--13 & FPGA & 1~GH/s      & Good \\
2013+    & ASIC & $>$100~TH/s & Best \\
\bottomrule
\end{tabular}
\end{table}

A modern Bitmain Antminer S19~XP delivers $\approx$140~TH/s at 3010~W.

\section{Mining Pools}

Solo mining variance is very high (expected time per solo miner: years).
\textbf{Mining pools} aggregate hash power:

\begin{itemize}
  \item Pool sends members \emph{share targets} (easier than block target).
  \item Members submit shares as proof of work.
  \item When pool finds a block, reward is split by shares.
  \item Common payout schemes: PPS (Pay-Per-Share), PPLNS.
\end{itemize}

\section{Why SHA256d and Not SHA256?}

Bitcoin uses \textbf{double SHA-256} (\texttt{SHA256d}) for block
hashing:
\[
  \text{block hash} = \mathrm{SHA256}(\mathrm{SHA256}(\text{header}))
\]

\textbf{Reason 1 — Length extension resistance}: SHA-256 has a length
extension vulnerability.  Applying it twice breaks the Merkle-Damg\aa rd
chain.

\textbf{Reason 2 — Extra security margin}: Even if SHA-256 were reduced
to a 128-bit effective security level, SHA256d provides an additional
layer.

\subsection{Nonce Space Exhaustion}

The 4-byte \texttt{nNonce} field gives only $2^{32}\approx 4$ billion
values.  At modern hashrates ($> 10^{18}$~H/s), a miner exhausts the
entire nonce space in microseconds.

\textbf{Solution}: Vary the \textbf{extra nonce} in the coinbase
transaction (first input's \texttt{scriptSig} can hold up to 100 bytes).
Changing the extra nonce changes the coinbase transaction, which changes
the Merkle root, which changes the block header — giving the miner a
fresh set of $2^{32}$ nonces to try.

In practice: modern mining hardware also varies the \texttt{nTime}
field (within $\pm 2$ hours of real time) for even more entropy.

\section{Selfish Mining}

A miner with $> 1/3$ of hashrate can mine a private chain, release it
strategically, and earn revenue disproportionate to their share.

\section{Transaction Lifecycle in the Mempool}

\begin{enumerate}
  \item User broadcasts transaction; nodes validate and add to \textbf{mempool}.
  \item Miner selects transactions (highest fee/vbyte first) into a block.
  \item Miner solves PoW; broadcasts block.
  \item After $\approx$6 confirmations ($\approx$60 min) the transaction
        is considered irreversible.
\end{enumerate}

\section{The Lightning Network}

Second-layer protocol built on Bitcoin for fast, cheap payments:
\begin{enumerate}
  \item Two parties open a \textbf{payment channel} — a 2-of-2 multisig UTXO.
  \item They exchange signed off-chain transactions (channel updates).
  \item Either party can close by broadcasting the latest signed state.
  \item Payments route through a network of channels (no direct channel needed).
\end{enumerate}

\section{Proof-of-Work vs.\ Alternatives}

\begin{table}[h]
\centering
\caption{Consensus mechanism comparison}
\begin{tabular}{p{2.5cm}p{3cm}p{3cm}p{3cm}}
\toprule
\textbf{Mechanism} & \textbf{Resource} & \textbf{Attack threshold} & \textbf{Examples} \\
\midrule
Proof of Work & Compute / energy & 50\% hashrate & Bitcoin \\
Proof of Stake & Staked coins & 33\% of stake (BFT) & Ethereum, Cardano \\
DPoS & Elected delegates & Delegate collusion & EOS, TRON \\
Proof of Space & Disk storage & 51\% storage & Chia \\
\bottomrule
\end{tabular}
\end{table}


% ================================================================
\chapter{Introduction to Ethereum}
\label{chap:eth_intro}
% ================================================================

\section{Ethereum vs.\ Bitcoin}

\begin{table}[h]
\centering
\caption{Bitcoin vs.\ Ethereum}
\begin{tabular}{lll}
\toprule
\textbf{Feature} & \textbf{Bitcoin} & \textbf{Ethereum} \\
\midrule
Primary purpose & Digital currency & Programmable platform \\
State model & UTXO & Account-based \\
Scripting & Limited (Script) & Turing-complete (EVM) \\
Block time & $\approx$10 min & 12 s (post-merge) \\
Consensus & PoW & PoS (since Sept 2022) \\
Supply cap & 21M BTC & None (issuance $\to$ 0) \\
Creator & Satoshi Nakamoto (2008) & Vitalik Buterin (2015) \\
\bottomrule
\end{tabular}
\end{table}

\section{Account Model}

\subsection{Externally Owned Accounts (EOAs)}

\begin{itemize}
  \item Controlled by a private key (human user).
  \item Can initiate transactions.
  \item Fields: \texttt{nonce}, \texttt{balance}.
  \item No code or storage.
\end{itemize}

\subsection{Contract Accounts}

\begin{itemize}
  \item Controlled by immutable EVM code.
  \item Created by a deployment transaction.
  \item Fields: \texttt{nonce}, \texttt{balance}, \texttt{storageRoot},
        \texttt{codeHash}.
  \item Executes when called by EOA or another contract.
\end{itemize}

\begin{align}
  \text{Account state} = (\underbrace{\texttt{nonce}}_{\text{tx count}},\;
  \underbrace{\texttt{balance}}_{\text{wei}},\;
  \underbrace{\texttt{storageRoot}}_{\text{trie hash}},\;
  \underbrace{\texttt{codeHash}}_{\text{keccak256(code)}})
\end{align}

\section{Ether and Wei}

\begin{align}
  1\;\text{ETH} &= 10^{18}\;\text{wei} \\
  1\;\text{ETH} &= 10^{9}\;\text{gwei}
\end{align}

Gas prices are quoted in \textbf{gwei}.

\section{Gas and Fees}

Gas is the unit of computational work inside the EVM.

\subsection{Why Gas?}

\begin{itemize}
  \item Prevents infinite loops (Turing halting problem) — execution stops
        when gas exhausted.
  \item Compensates validators for computation.
  \item Creates a market for block space.
\end{itemize}

\subsection{EIP-1559 Fee Mechanism (Post-London 2021)}

\begin{align}
  \text{effective gas price} &= \min(\texttt{maxFeePerGas},\;
    \texttt{baseFee}+\texttt{priorityFee}) \\
  \text{fee paid} &= \text{gasUsed}\times\text{effective gas price}
\end{align}

\begin{itemize}
  \item \textbf{Base fee}: Algorithmically set; \emph{burned} each block.
  \item \textbf{Priority fee (tip)}: Given to the validator.
  \item Base fee rises 12.5\% if previous block was full; falls 12.5\% if
        empty (target: 50\% full).
\end{itemize}

\subsection{EIP-1559: Worked Fee Example}

\begin{example}
  \textbf{Given}: base fee = 20~gwei, user sets
  \texttt{maxPriorityFeePerGas} = 2~gwei,
  \texttt{maxFeePerGas} = 25~gwei.  Transaction uses 21\,000 gas
  (simple ETH transfer).

  \textbf{Step 1}: Effective gas price
  = $\min(25, 20+2) = 22$~gwei.

  \textbf{Step 2}: Fee paid = $21000\times22 = 462000$~gwei
  $= 0.000462$~ETH.

  \textbf{Step 3}: Base fee burned = $21000\times20 = 420000$~gwei.

  \textbf{Step 4}: Validator tip = $21000\times2 = 42000$~gwei.

  \textbf{Step 5}: Refund to user = $(25-22)\times21000 = 63000$~gwei.

  \textbf{Summary}: User paid 0.000462~ETH; 0.000420~ETH was burned;
  validator received 0.000042~ETH.
\end{example}

\begin{keyinsight}
  \textbf{Deflationary effect}: Base fee is \emph{burned} (removed from
  circulation) rather than given to validators.  During periods of high
  activity, ETH supply decreases.  In 2021--2022 Ethereum was briefly
  \emph{deflationary} (more ETH burned than issued).
\end{keyinsight}

\subsection{Account Nonce and Replay Protection}

Every transaction from an EOA includes a \textbf{nonce} equal to the
number of transactions previously sent from that address.

\begin{itemize}
  \item Prevents \textbf{replay attacks}: if Alice sends a transaction,
        an attacker cannot replay the same transaction (the nonce would
        be wrong).
  \item Enforces \textbf{transaction ordering}: transactions from the
        same sender are executed in nonce order.
  \item If a nonce is skipped (e.g., nonce 3 but nonce 2 not yet
        mined), the transaction is held in the mempool until its
        predecessor arrives.
\end{itemize}

\textbf{EIP-155 replay protection}: The transaction includes
\texttt{chainId} in its signature hash.  A transaction signed for
Ethereum mainnet (chainId=1) cannot be replayed on Ethereum Classic
(chainId=61) or any other chain.

\section{Smart Contracts}

\begin{definitionbox}
  A \textbf{smart contract} is a program deployed on the blockchain that
  executes deterministically on the EVM, has its own address/balance/
  storage, is immutable once deployed, and is triggered by transactions.
\end{definitionbox}

\subsection{Representative Applications}

\begin{description}
  \item[DeFi] Uniswap (DEX), Aave/Compound (lending), MakerDAO (stablecoins).
  \item[NFTs] ERC-721 unique digital assets.
  \item[DAOs] On-chain governance (voting, treasury management).
  \item[Privacy] Tornado Cash (Chapter~\ref{chap:tornado}).
\end{description}

\section{ERC-20 Token Standard}

Standard interface for fungible tokens:
\begin{itemize}
  \item \texttt{totalSupply()} — total token supply.
  \item \texttt{balanceOf(address)} — balance of an address.
  \item \texttt{transfer(to, amount)} — transfer tokens.
  \item \texttt{approve(spender, amount)} — allow delegation.
  \item \texttt{transferFrom(from, to, amount)} — spend delegated tokens.
  \item Events: \texttt{Transfer}, \texttt{Approval}.
\end{itemize}


% ================================================================
\chapter{Ethereum Data Structures and Encoding}
\label{chap:ethdata}
% ================================================================

\section{RLP Encoding}

\textbf{Recursive Length Prefix (RLP)} is the canonical serialisation
format in Ethereum for all data structures.

\subsection{Encoding Rules}

\textbf{Single byte $b$}:
\begin{itemize}
  \item $b\in[0x00,\,0x7f]$: encode as $b$ itself.
\end{itemize}

\textbf{Byte string of length $\ell$}:
\begin{itemize}
  \item $\ell=0$: \texttt{0x80}.
  \item $\ell\leq 55$: $(0x80+\ell)\;\|\;$string.
  \item $\ell>55$: $(0xb7+\text{len}(\ell))\;\|\;\text{BE}(\ell)\;\|\;$string.
\end{itemize}

\textbf{List with total payload length $\ell$}:
\begin{itemize}
  \item $\ell\leq 55$: $(0xc0+\ell)\;\|\;$RLP-encoded items.
  \item $\ell>55$: $(0xf7+\text{len}(\ell))\;\|\;\text{BE}(\ell)\;\|\;$items.
\end{itemize}

\subsection{Examples}

\begin{itemize}
  \item \texttt{RLP("dog")} $=$ \texttt{0x83 64 6f 67}
        \quad ($\ell=3$, $0x80+3=0x83$)
  \item \texttt{RLP(["cat","dog"])} $=$ \texttt{0xc8 83 63 61 74 83 64 6f 67}
        \quad (payload $=8$ bytes, $0xc0+8=0xc8$)
  \item \texttt{RLP("")} $=$ \texttt{0x80}
  \item \texttt{RLP([])} $=$ \texttt{0xc0}
\end{itemize}

\section{Merkle Patricia Trie (MPT)}

Ethereum stores its state in a \textbf{Merkle Patricia Trie}: a
compressed radix trie where every node is identified by its hash.

\subsection{Node Types}

\begin{description}
  \item[Empty node] Null — represents missing key.
  \item[Leaf node] $(k_{\text{enc}},\;\text{value})$ — stores a
        key-value pair.
  \item[Extension node] $(k_{\text{enc}},\;\text{hash\_of\_next})$ —
        compresses a shared key prefix.
  \item[Branch node] 17 children: indices 0--f (hex nibble) plus an
        optional value slot.
\end{description}

\subsection{Hex-Prefix (HP) Path Encoding}

Keys are sequences of hex nibbles; a flag byte encodes type and parity:
\begin{table}[h]
\centering
\begin{tabular}{cll}
\toprule
\textbf{Prefix nibble} & \textbf{Node type} & \textbf{Path length} \\
\midrule
\texttt{0} & Extension & Even \\
\texttt{1} & Extension & Odd \\
\texttt{2} & Leaf      & Even \\
\texttt{3} & Leaf      & Odd \\
\bottomrule
\end{tabular}
\end{table}

\section{MPT Worked Example}

\begin{example}
  Build a Merkle Patricia Trie with two key-value pairs:
  \begin{itemize}
    \item Key \texttt{0x1234} $\to$ value "dog"
    \item Key \texttt{0x1235} $\to$ value "cat"
  \end{itemize}

  Both keys share the nibble prefix \texttt{1,2,3}.

  \textbf{Step 1}: Both paths share prefix 1,2,3.  Create an
  \textbf{extension node} with key \texttt{1,2,3} pointing to a branch.

  \textbf{Step 2}: After the shared prefix, the paths diverge at nibble
  4 vs 5.  Create a \textbf{branch node} with children at positions 4
  and 5.

  \textbf{Step 3}: Each child is a \textbf{leaf node}:
  \begin{itemize}
    \item Position 4: leaf (remaining path = empty, value = "dog")
    \item Position 5: leaf (remaining path = empty, value = "cat")
  \end{itemize}

  \textbf{Resulting structure}:
\begin{lstlisting}[language={}]
Extension(path=[1,2,3], next=Branch)
  Branch[4] -> Leaf(path=[], value="dog")
  Branch[5] -> Leaf(path=[], value="cat")
\end{lstlisting}

  \textbf{Merkle hashing}: Each node is hashed with keccak256.  The root
  hash is the hash of the extension node.  Changing ``dog'' to ``duck''
  changes the leaf hash, the branch hash, and the extension (root) hash.
  This is the Merkle property.
\end{example}

\section{The Four Ethereum Tries}

\begin{table}[h]
\centering
\caption{Tries referenced in block headers}
\begin{tabular}{p{3cm}p{4cm}p{5cm}}
\toprule
\textbf{Trie} & \textbf{Key} & \textbf{Value (RLP-encoded)} \\
\midrule
State          & \texttt{keccak256(address)} & (nonce, balance, storageRoot, codeHash) \\
Transactions   & \texttt{RLP(tx index)} & serialised transaction \\
Receipts       & \texttt{RLP(tx index)} & (status, gasUsed, logs, bloom) \\
Storage (per contract) & \texttt{keccak256(slot)} & 256-bit value \\
\bottomrule
\end{tabular}
\end{table}

\begin{keyinsight}
  The \texttt{stateRoot} in the block header is the root hash of the
  global state trie.  Changing any account (balance, nonce, storage)
  changes the root, making tampering detectable.
\end{keyinsight}

\section{Keccak-256}

Ethereum uses \textbf{Keccak-256} (a variant of SHA-3) for:
\begin{itemize}
  \item Address derivation: last 20 bytes of keccak256(uncompressed pubkey).
  \item EVM \texttt{KECCAK256} opcode.
  \item Function selectors: first 4 bytes of keccak256(signature string).
  \item All MPT node hashes.
\end{itemize}

\section{Bloom Filters}

Each block header contains a 2048-bit \textbf{Bloom filter}:
\begin{itemize}
  \item Compact representation of all addresses and topics in the block's logs.
  \item False positives possible; false negatives impossible.
  \item Enables efficient log searching without a full node.
\end{itemize}


% ================================================================
\chapter{Ethereum Transactions}
\label{chap:ethtx}
% ================================================================

\section{Transaction Types}

\begin{description}
  \item[Type~0 (Legacy)] Single \texttt{gasPrice} field; pre-EIP-1559.
  \item[Type~1 (EIP-2930)] Adds \texttt{accessList} for storage slot pre-warming.
  \item[Type~2 (EIP-1559)] \texttt{maxFeePerGas} + \texttt{maxPriorityFeePerGas};
        dominant format today.
\end{description}

\section{EIP-1559 Transaction Fields}

\begin{table}[h]
\centering
\caption{Type-2 transaction fields}
\begin{tabular}{lll}
\toprule
\textbf{Field} & \textbf{Type} & \textbf{Description} \\
\midrule
\texttt{chainId}              & uint256 & Prevents cross-chain replay \\
\texttt{nonce}                & uint256 & Sequential tx count from sender \\
\texttt{maxPriorityFeePerGas} & uint256 & Max tip to validator (wei/gas) \\
\texttt{maxFeePerGas}         & uint256 & Max total fee (wei/gas) \\
\texttt{gasLimit}             & uint256 & Max gas units allowed \\
\texttt{to}                   & address & Recipient (\texttt{null} = deploy) \\
\texttt{value}                & uint256 & ETH to transfer (wei) \\
\texttt{data}                 & bytes   & Calldata / initcode \\
\texttt{accessList}           & list    & Pre-warmed slots \\
\texttt{v, r, s}              &         & ECDSA signature \\
\bottomrule
\end{tabular}
\end{table}

\section{Signing and Sender Recovery}

No explicit \texttt{from} field; the sender is recovered from the signature:
\begin{align}
  \texttt{sigHash} &= \mathrm{keccak256}(\mathrm{RLP}(2,\,\texttt{chainId},\,\texttt{nonce},\ldots,\texttt{accessList})) \\
  (v,r,s) &= \mathrm{ECDSA\_sign}(\texttt{sigHash},\;\text{private key}) \\
  \texttt{sender} &= \mathrm{ecrecover}(v,r,s,\;\texttt{sigHash})
\end{align}

\section{Transaction Execution Flow}

\begin{enumerate}
  \item Validate: \texttt{nonce} matches, sender balance $\geq$ \texttt{gasLimit}$\times$\texttt{maxFeePerGas}.
  \item Deduct \texttt{gasLimit}$\times$\texttt{maxFeePerGas} from sender upfront.
  \item Execute EVM code (or simple value transfer).
  \item Refund unused gas: $(\texttt{gasLimit}-\texttt{gasUsed})\times\texttt{effectivePrice}$.
  \item \textbf{Burn} base fee: $\texttt{gasUsed}\times\texttt{baseFee}$.
  \item Pay validator: $\texttt{gasUsed}\times\texttt{priorityFee}$.
\end{enumerate}

\section{Contract Creation}

When \texttt{to} is null:
\begin{itemize}
  \item \texttt{data} = \textbf{initcode} (constructor bytecode + arguments).
  \item EVM runs initcode; constructor returns runtime bytecode.
  \item Contract address (CREATE):
        $\texttt{addr}=\mathrm{keccak256}(\mathrm{RLP}(\texttt{sender},\texttt{nonce}))[12\text{:}]$
  \item Contract address (CREATE2):
        $\texttt{addr}=\mathrm{keccak256}(\texttt{0xff}\|\texttt{sender}\|\texttt{salt}\|\mathrm{keccak256(initcode)})[12\text{:}]$
\end{itemize}

\section{Function Calls}

\texttt{data} = 4-byte selector $\|$ ABI-encoded arguments.

\begin{align}
  \texttt{selector} = \mathrm{keccak256}(\text{"funcName(type1,type2,...)")}[0:4]
\end{align}

Example — calling \texttt{transfer(address,uint256)}:
\begin{lstlisting}[language={}]
selector  = 0xa9059cbb
calldata  = selector || abi.encode(to, amount)
\end{lstlisting}

\section{Internal Calls}

\begin{table}[h]
\centering
\caption{EVM call types}
\begin{tabular}{lll}
\toprule
\textbf{Opcode} & \textbf{Context} & \textbf{State changes} \\
\midrule
\texttt{CALL}         & Callee's context & Yes \\
\texttt{STATICCALL}   & Callee's context & No (read-only) \\
\texttt{DELEGATECALL} & Caller's context & Yes (in caller's storage) \\
\bottomrule
\end{tabular}
\end{table}


% ================================================================
\chapter{Ethereum Virtual Machine (EVM)}
\label{chap:evm}
% ================================================================

\section{Architecture}

The EVM is a stack-based, 256-bit quasi-Turing-complete virtual machine:

\begin{itemize}
  \item \textbf{Stack}: LIFO, max 1024 elements, each 256 bits.
  \item \textbf{Memory}: Volatile byte array, expands on demand (cost $\propto$ size$^2$ for large allocations).
  \item \textbf{Storage}: Persistent 256-bit $\to$ 256-bit key-value map, per contract.
  \item \textbf{Calldata}: Read-only input from the calling transaction.
  \item \textbf{Code}: Read-only bytecode being executed.
\end{itemize}

\section{Instruction Set}

\begin{table}[h]
\centering
\caption{Key EVM opcodes and gas costs}
\begin{tabular}{lllp{4cm}}
\toprule
\textbf{Category} & \textbf{Opcode} & \textbf{Gas} & \textbf{Effect} \\
\midrule
Arithmetic & \texttt{ADD} & 3 & $a+b$ \\
 & \texttt{MUL} & 5 & $a\times b$ \\
 & \texttt{EXP} & 10+ & $a^b$ \\
\midrule
Bitwise & \texttt{AND}, \texttt{OR}, \texttt{XOR} & 3 & bitwise ops \\
 & \texttt{SHL}, \texttt{SHR} & 3 & shifts \\
\midrule
Memory & \texttt{MLOAD} & 3+ & load 32 B \\
 & \texttt{MSTORE} & 3+ & store 32 B \\
\midrule
Storage & \texttt{SLOAD} & 2100 & load slot \\
 & \texttt{SSTORE} & 20\,000 & write new slot \\
\midrule
Control & \texttt{JUMP}, \texttt{JUMPI} & 8/10 & unconditional/conditional \\
 & \texttt{REVERT} & 0 & revert + return data \\
 & \texttt{RETURN} & 0 & return data \\
\midrule
Context & \texttt{CALLER} & 2 & \texttt{msg.sender} \\
 & \texttt{CALLVALUE} & 2 & \texttt{msg.value} \\
 & \texttt{TIMESTAMP} & 2 & block timestamp \\
\midrule
Hash & \texttt{KECCAK256} & 30+6/word & keccak256(mem range) \\
\midrule
Calls & \texttt{CALL} & 100+ & external call \\
 & \texttt{STATICCALL} & 100+ & read-only call \\
 & \texttt{DELEGATECALL} & 100+ & delegate call \\
\bottomrule
\end{tabular}
\end{table}

\section{Execution Trace Example}

\begin{lstlisting}[language={},caption=Computing $3+5$ in EVM bytecode]
PUSH1 0x03   // Stack: [3]
PUSH1 0x05   // Stack: [5, 3]
ADD          // Stack: [8]
\end{lstlisting}

\section{Gas Accounting in Practice}

\begin{example}
  \textbf{Executing a simple ERC-20 transfer}:

  Approximate gas breakdown for \texttt{transfer(address,uint256)}:
  \begin{table}[h]
  \centering
  \begin{tabular}{lr}
  \toprule
  \textbf{Operation} & \textbf{Gas} \\
  \midrule
  Base transaction fee & 21\,000 \\
  Non-zero calldata (68 bytes $\times$ 16) & 1\,088 \\
  SLOAD (read sender balance) & 2\,100 \\
  SLOAD (read recipient balance) & 2\,100 \\
  SSTORE (update sender balance) & 5\,000 \\
  SSTORE (update recipient balance) & 20\,000 \\
  LOG3 (Transfer event) & 1\,500 \\
  Misc arithmetic/logic & $\approx 500$ \\
  \midrule
  \textbf{Total} & $\approx 53\,288$ \\
  \bottomrule
  \end{tabular}
  \end{table}
\end{example}

\subsection{Out-of-Gas and Reverts}

If gas is exhausted mid-execution:
\begin{itemize}
  \item An \textbf{out-of-gas} exception is raised.
  \item All state changes from this call (and any sub-calls) are
        \emph{reverted}.
  \item All remaining gas is consumed (none refunded).
  \item The transaction is included in the block (fee paid) but has
        status 0 (failed).
\end{itemize}

If \texttt{REVERT} is executed explicitly:
\begin{itemize}
  \item State changes are reverted.
  \item \emph{Unused} gas is refunded.
  \item A return data (error message) is available.
\end{itemize}

\subsection{Call Stack and Call Depth Limit}

When a contract calls another, a new \textbf{call frame} is pushed onto
the call stack:
\begin{itemize}
  \item Each frame has its own \textbf{stack}, \textbf{memory}, and
        \textbf{program counter}.
  \item Storage and balance are global (shared across frames).
  \item The call depth limit is \textbf{1024}.  Exceeding it causes an
        out-of-gas-like failure.
  \item The EIP-150 \textbf{63/64 rule}: a CALL forwards at most
        $\frac{63}{64}$ of remaining gas, ensuring gas eventually runs
        out even in deep recursion.
\end{itemize}

\section{ABI Encoding}

\textbf{Static types} (\texttt{uint256}, \texttt{address}, \texttt{bool}):
padded to 32 bytes, left-aligned.

\textbf{Dynamic types} (\texttt{bytes}, \texttt{string}, dynamic arrays):
encoded with a head containing an offset to the tail data.

\section{Contract Deployment Bytecode}

\begin{enumerate}
  \item \textbf{Constructor code}: copies runtime bytecode to memory, returns it.
  \item \textbf{Runtime bytecode}: stored at contract address; executed on calls.
  \item \textbf{Constructor arguments}: ABI-encoded, appended after runtime bytecode.
\end{enumerate}

\begin{exambox}
  \texttt{SSTORE} to a new (zero) slot costs 20\,000 gas; updating an
  existing slot costs 5\,000; clearing to zero gives a 15\,000 gas
  refund (capped at 20\% of gas used in EIP-3529).
\end{exambox}


% ================================================================
\chapter{Ethereum Blocks}
\label{chap:ethblocks}
% ================================================================

\section{Block Structure}

An Ethereum block = \textbf{execution payload} (transactions, state)
wrapped in a \textbf{beacon block} (consensus metadata).

\subsection{Execution Block Header}

\begin{table}[h]
\centering
\caption{Execution layer block header (post-merge)}
\begin{tabular}{lll}
\toprule
\textbf{Field} & \textbf{Type} & \textbf{Description} \\
\midrule
\texttt{parentHash}       & bytes32  & Hash of parent block header \\
\texttt{ommersHash}       & bytes32  & keccak256 of uncle list (always empty post-merge) \\
\texttt{beneficiary}      & address  & Fee recipient \\
\texttt{stateRoot}        & bytes32  & Root of state MPT after all txs \\
\texttt{transactionsRoot} & bytes32  & Root of tx trie \\
\texttt{receiptsRoot}     & bytes32  & Root of receipts trie \\
\texttt{logsBloom}        & bytes256 & Bloom filter for logs \\
\texttt{difficulty}       & uint256  & 0 post-merge \\
\texttt{number}           & uint256  & Block height \\
\texttt{gasLimit}         & uint256  & Max total gas \\
\texttt{gasUsed}          & uint256  & Total gas consumed \\
\texttt{timestamp}        & uint64   & Unix time (exact 12 s slots) \\
\texttt{extraData}        & bytes    & Up to 32 bytes arbitrary \\
\texttt{prevRandao}       & bytes32  & Beacon chain RANDAO mix \\
\texttt{nonce}            & bytes8   & 0 post-merge \\
\texttt{baseFeePerGas}    & uint256  & EIP-1559 base fee \\
\texttt{withdrawalsRoot}  & bytes32  & Validator withdrawals (EIP-4895) \\
\bottomrule
\end{tabular}
\end{table}

\section{Transaction Receipts}

\begin{itemize}
  \item \texttt{status}: 1 = success, 0 = revert.
  \item \texttt{cumulativeGasUsed}: running total in block.
  \item \texttt{logs}: events emitted.
  \item \texttt{logsBloom}: Bloom contribution.
\end{itemize}

\section{Logs and Events}

\begin{lstlisting}[language=Solidity]
event Transfer(address indexed from, address indexed to, uint256 value);
\end{lstlisting}

Log structure:
\begin{itemize}
  \item Topic 0: \texttt{keccak256("Transfer(address,address,uint256)")}
  \item Topic 1: \texttt{from} (indexed)
  \item Topic 2: \texttt{to} (indexed)
  \item Data: ABI-encoded \texttt{value}
\end{itemize}

Events are \emph{not} accessible from within the EVM; they are only
visible to external observers.

\section{The Merge (September 15, 2022)}

\begin{itemize}
  \item Energy reduced 99.95\%: 5.13~GW $\to$ 2.62~MW.
  \item PoW nonce/difficulty set to 0.
  \item \texttt{prevRandao} replaces \texttt{mixHash}, injecting beacon
        chain randomness.
  \item Block time became exactly 12 s (one slot).
\end{itemize}


% ================================================================
\chapter{Solidity Programming}
\label{chap:solidity}
% ================================================================

\section{Language Overview}

Solidity is a statically-typed, object-oriented language that compiles
to EVM bytecode.  It supports contracts, inheritance, events, and
custom modifiers.

\begin{lstlisting}[language=Solidity,caption=Minimal contract]
// SPDX-License-Identifier: MIT
pragma solidity ^0.8.0;

contract SimpleStorage {
    uint256 private value;

    function set(uint256 _v) public { value = _v; }

    function get() public view returns (uint256) { return value; }
}
\end{lstlisting}

\section{Types}

\subsection{Value Types}
\texttt{uint8}--\texttt{uint256}, \texttt{int8}--\texttt{int256},
\texttt{bool}, \texttt{address} / \texttt{address payable},
\texttt{bytes1}--\texttt{bytes32}, \texttt{enum}.

\subsection{Reference Types}
\texttt{bytes} (dynamic), \texttt{string}, arrays (\texttt{T[]}),
\texttt{struct}, \texttt{mapping(K => V)}.

\subsection{Special Variables}
\begin{itemize}
  \item \texttt{msg.sender} — caller address.
  \item \texttt{msg.value} — ETH sent (wei).
  \item \texttt{block.timestamp}, \texttt{block.number}.
  \item \texttt{tx.origin} — original EOA (avoid for auth).
\end{itemize}

\section{Function Visibility and Mutability}

\begin{table}[h]
\centering
\caption{Function visibility modifiers}
\begin{tabular}{ll}
\toprule
\textbf{Visibility} & \textbf{Accessible from} \\
\midrule
\texttt{public}   & Anywhere (EOAs, other contracts, internally) \\
\texttt{external} & Only from outside (more efficient calldata access) \\
\texttt{internal} & Within contract and derived contracts \\
\texttt{private}  & Only within this contract \\
\bottomrule
\end{tabular}
\end{table}

State mutability: \texttt{view} (reads state), \texttt{pure} (no state
access), \texttt{payable} (can receive ETH), default (reads + writes).

\section{Mappings, Structs, Events}

\begin{lstlisting}[language=Solidity,caption=Bank contract with events]
contract Bank {
    mapping(address => uint256) public balances;
    event Deposit(address indexed user, uint256 amount);
    event Withdrawal(address indexed user, uint256 amount);

    function deposit() public payable {
        balances[msg.sender] += msg.value;
        emit Deposit(msg.sender, msg.value);
    }

    function withdraw(uint256 amount) public {
        require(balances[msg.sender] >= amount, "Insufficient");
        balances[msg.sender] -= amount;       // effect first!
        (bool ok, ) = msg.sender.call{value: amount}("");
        require(ok, "Transfer failed");
        emit Withdrawal(msg.sender, amount);
    }
}
\end{lstlisting}

\begin{exambox}[Reentrancy]
  Always update state \emph{before} external calls.  The DAO hack (2016)
  drained \$60~M by re-entering a withdraw function before the balance
  was updated.  Pattern: \textbf{Checks $\to$ Effects $\to$ Interactions}.
\end{exambox}

\section{Modifiers and Inheritance}

\begin{lstlisting}[language=Solidity,caption=Ownable pattern]
contract Ownable {
    address public owner;
    constructor() { owner = msg.sender; }

    modifier onlyOwner() {
        require(msg.sender == owner, "Not owner");
        _;
    }
}

contract Token is Ownable {
    mapping(address => uint256) public balanceOf;

    function mint(address to, uint256 amount) public onlyOwner {
        balanceOf[to] += amount;
    }
}
\end{lstlisting}

\section{Minimal ERC-20 Implementation}

\begin{lstlisting}[language=Solidity,caption=ERC-20 token]
// SPDX-License-Identifier: MIT
pragma solidity ^0.8.0;

contract ERC20 {
    string  public name;
    string  public symbol;
    uint8   public decimals = 18;
    uint256 public totalSupply;

    mapping(address => uint256) public balanceOf;
    mapping(address => mapping(address => uint256)) public allowance;

    event Transfer(address indexed from, address indexed to, uint256 value);
    event Approval(address indexed owner, address indexed spender, uint256 value);

    constructor(string memory _name, string memory _symbol, uint256 _supply) {
        name = _name; symbol = _symbol;
        totalSupply = _supply * 10**decimals;
        balanceOf[msg.sender] = totalSupply;
    }

    function transfer(address to, uint256 v) public returns (bool) {
        require(balanceOf[msg.sender] >= v);
        balanceOf[msg.sender] -= v;
        balanceOf[to] += v;
        emit Transfer(msg.sender, to, v);
        return true;
    }

    function approve(address spender, uint256 v) public returns (bool) {
        allowance[msg.sender][spender] = v;
        emit Approval(msg.sender, spender, v);
        return true;
    }

    function transferFrom(address from, address to, uint256 v)
        public returns (bool)
    {
        require(balanceOf[from] >= v && allowance[from][msg.sender] >= v);
        allowance[from][msg.sender] -= v;
        balanceOf[from] -= v;
        balanceOf[to]   += v;
        emit Transfer(from, to, v);
        return true;
    }
}
\end{lstlisting}

\section{Storage Layout}

State variables are stored in 32-byte \textbf{slots} starting at slot 0:
\begin{itemize}
  \item Variables are packed into the same slot if they fit.
  \item Mappings: \texttt{mapping(K=>V)} stores value at
        \texttt{keccak256(K\,||\,slot)}.
  \item Dynamic arrays: length at slot $s$; elements at
        \texttt{keccak256(s) + i}.
\end{itemize}

\begin{lstlisting}[language=Solidity,caption=Storage layout example]
contract Layout {
    uint128 a;    // slot 0, bytes 0-15
    uint128 b;    // slot 0, bytes 16-31  (packed with a)
    uint256 c;    // slot 1               (full slot)
    address d;    // slot 2, bytes 0-19
    bool    e;    // slot 2, byte 20      (packed with d)
    mapping(address => uint256) balances; // slot 3 (base)
    // balances[addr] is at keccak256(abi.encode(addr, 3))
}
\end{lstlisting}

Understanding storage layout is critical for:
\begin{itemize}
  \item Writing \texttt{delegatecall}-based proxies (storage must match).
  \item Optimising gas (packing saves SSTORE costs).
  \item Security auditing (spotting storage collisions).
\end{itemize}

\section{Interfaces and Abstract Contracts}

\begin{lstlisting}[language=Solidity,caption=Interface definition]
interface IERC20 {
    function totalSupply() external view returns (uint256);
    function balanceOf(address) external view returns (uint256);
    function transfer(address to, uint256 amount) external returns (bool);
    function approve(address spender, uint256 amount) external returns (bool);
    function transferFrom(address from, address to, uint256 amount)
        external returns (bool);
    event Transfer(address indexed from, address indexed to, uint256 value);
    event Approval(address indexed owner, address indexed spender, uint256 value);
}
\end{lstlisting}

Interfaces allow contracts to interact with any ERC-20 token without
knowing its implementation:
\begin{lstlisting}[language=Solidity]
IERC20 token = IERC20(tokenAddress);
uint256 bal = token.balanceOf(msg.sender);
\end{lstlisting}

\section{The Proxy Pattern (Upgradeable Contracts)}

Smart contracts are immutable.  The proxy pattern separates logic from
storage to allow upgrades:

\begin{enumerate}
  \item A \textbf{proxy contract} holds all storage and ETH.
  \item It \texttt{delegatecall}s a \textbf{logic contract} for all
        function calls.
  \item To upgrade: deploy a new logic contract and update the proxy's
        stored implementation address.
\end{enumerate}

\begin{lstlisting}[language=Solidity,caption=Minimal transparent proxy]
contract Proxy {
    address public implementation;   // slot 0
    address public admin;            // slot 1

    constructor(address _impl) {
        implementation = _impl;
        admin = msg.sender;
    }

    function upgrade(address newImpl) external {
        require(msg.sender == admin);
        implementation = newImpl;
    }

    fallback() external payable {
        address impl = implementation;
        assembly {
            calldatacopy(0, 0, calldatasize())
            let result := delegatecall(gas(), impl, 0, calldatasize(), 0, 0)
            returndatacopy(0, 0, returndatasize())
            switch result
            case 0 { revert(0, returndatasize()) }
            default { return(0, returndatasize()) }
        }
    }
}
\end{lstlisting}

\begin{exambox}
  \textbf{Storage collision hazard}: If the proxy stores
  \texttt{implementation} at slot 0 and the logic contract also uses
  slot 0 for a state variable, a call will overwrite the implementation
  pointer!  EIP-1967 standardises proxy storage slots to
  \texttt{keccak256("eip1967.proxy.implementation") - 1} to avoid this.
\end{exambox}

\section{Common Vulnerabilities}

\begin{table}[h]
\centering
\caption{Common Solidity vulnerabilities}
\begin{tabular}{p{3cm}p{4.5cm}p{4cm}}
\toprule
\textbf{Vulnerability} & \textbf{Description} & \textbf{Mitigation} \\
\midrule
Reentrancy & Re-enter before state updated & CEI pattern; \texttt{nonReentrant} \\
Integer overflow & Wrap-around (pre-0.8) & Solidity $\geq 0.8$; SafeMath \\
Missing access ctrl & No \texttt{onlyOwner} & OpenZeppelin AccessControl \\
Front-running & Mempool visibility & Commit-reveal \\
\texttt{tx.origin} auth & Phishing via forwarding & Always use \texttt{msg.sender} \\
\texttt{delegatecall} & Storage layout clash & Careful proxy patterns \\
\bottomrule
\end{tabular}
\end{table}


% ================================================================
\chapter{JavaScript and Node.js}
\label{chap:js}
% ================================================================

\section{JavaScript Essentials}

\begin{lstlisting}[language=JavaScript,caption=Variables and functions]
let x = 1;          // block-scoped, reassignable
const PI = 3.14;    // block-scoped, constant
var y = 2;          // function-scoped -- avoid

// Arrow function
const add = (a, b) => a + b;

// Arrow with block body
const safeDivide = (a, b) => {
    if (b === 0) throw new Error("Division by zero");
    return a / b;
};
\end{lstlisting}

\section{Arrays}

\begin{lstlisting}[language=JavaScript,caption=Array operations]
let arr = [1, 2, 'Hello'];
arr.push(4);          // Add to end
arr.pop();            // Remove from end
arr.unshift('A');     // Add to start
arr.shift();          // Remove from start

// Functional methods
arr.map(x => x * 2);
arr.filter(x => typeof x === 'number');
arr.reduce((acc, x) => acc + x, 0);
\end{lstlisting}

\section{Objects}

\begin{lstlisting}[language=JavaScript,caption=Objects]
let user = { name: "John", age: 30 };

console.log(user.name);      // dot notation
console.log(user["age"]);    // bracket notation

user.isAdmin = true;         // add property
delete user.age;             // remove property
console.log("age" in user);  // false

for (let k in user) { console.log(k, user[k]); }
\end{lstlisting}

\section{Promises and Async/Await}

Many browser and Node.js APIs are asynchronous.

\begin{lstlisting}[language=JavaScript,caption=Promise chain]
fetch("https://example.com/")
    .then(response => response.text())
    .then(text => console.log(text))
    .catch(e => console.log("Error:", e));
\end{lstlisting}

\texttt{async}/\texttt{await} is syntactic sugar making async code look
synchronous:

\begin{lstlisting}[language=JavaScript,caption=async/await]
async function getText(url) {
    const response = await fetch(url);
    const text     = await response.text();
    console.log(text);
}
getText("https://example.com/");
\end{lstlisting}

An \texttt{async} function always returns a \texttt{Promise}.
\texttt{await} can only be used inside \texttt{async} functions.

\section{Node.js}

\textbf{Node.js} enables JavaScript outside the browser:
\begin{itemize}
  \item Built on Chrome's V8 engine.
  \item Event-driven, non-blocking I/O.
  \item \textbf{npm} (Node Package Manager): millions of packages.
\end{itemize}

\begin{lstlisting}[language=bash,caption=npm essentials]
npm init                     # create package.json
npm install ethers           # install locally
npm install -g truffle       # install globally
node script.js               # run a script
\end{lstlisting}

\section{Accepting CLI Input with Inquirer}

\begin{lstlisting}[language=JavaScript,caption=inquirer example (main.js)]
import inquirer from "inquirer";   // npm install inquirer

const questions = [
    { name: "name",    type: "input",
      message: "What's your name?" },
    { name: "program", type: "select",
      choices: ["BTech","Dual Degree","MTech","IDDP","PhD"],
      message: "Which program are you in?" },
];

inquirer.prompt(questions).then(answers => {
    console.log(`Name: ${answers.name}`);
    console.log(`Program: ${answers.program}`);
});
\end{lstlisting}

\texttt{package.json} must include \texttt{"type":"module"} to use
\texttt{import} syntax.

\section{Connecting to Ethereum with ethers.js}

\begin{lstlisting}[language=JavaScript,caption=ethers.js basics]
import { ethers } from "ethers";

const provider = new ethers.BrowserProvider(window.ethereum);
const signer   = await provider.getSigner();

const contract = new ethers.Contract(
    contractAddress, contractABI, signer);

const balance = await contract.balanceOf(signer.address);

const tx = await contract.transfer(recipient, amount);
await tx.wait();   // wait for 1 confirmation
\end{lstlisting}


% ================================================================
\chapter{Tornado Cash and Privacy}
\label{chap:tornado}
% ================================================================

\section{Motivation}

All Ethereum transactions are public.  Any observer who knows your
address can track every deposit and withdrawal, revealing your total
holdings.

\textbf{Concrete risk}: Pay a hotel 1~ETH; the front-desk clerk sees
your wallet holds 99~ETH and has kidnapping connections.

\textbf{Option A} — use an exchange: custody risk, KYC leaks.\\
\textbf{Option B} — send to a fresh address first: repeated
transfers still link old and new address.\\
\textbf{Tornado Cash} is a better Option B: a smart-contract
\emph{mixer} that breaks the on-chain link.

\section{Protocol Overview}

\begin{enumerate}
  \item Alice, Bob, Carol each deposit 1~ETH into the Tornado Cash
        contract.
  \item Later, anyone who holds a valid private note can withdraw 1~ETH
        to any fresh address.
  \item On-chain, the withdrawal cannot be linked to any specific deposit.
\end{enumerate}

\subsection{Desired Properties}

\begin{description}
  \item[Soundness] Only a genuine depositor can withdraw; no double withdrawal.
  \item[Privacy]   A withdrawal is unlinkable to its deposit.
\end{description}

\section{Commitment Scheme}

\subsection{Private Note}

The user's browser generates 62 random bytes:
\begin{lstlisting}[language={}]
tornado-eth-0.1-5-0x<Header><31-byte Nullifier><31-byte Secret>
\end{lstlisting}

\subsection{Pedersen Hash}

A \textbf{commitment} is the Pedersen hash of the concatenated bytes:
\begin{equation}
  C = \mathrm{PedersenHash}(b_1 b_2\cdots b_n)
    = g_1^{b_1} g_2^{b_2}\cdots g_n^{b_n}
\end{equation}
where $g_i$ are fixed public generators of an elliptic curve group.
The commitment is \emph{hiding} (reveals nothing about nullifier/secret)
and \emph{binding} (cannot find two inputs with the same hash).

\section{Deposit Steps}

\begin{enumerate}
  \item Generate random nullifier and secret on client.
  \item Compute $C = \mathrm{PedersenHash}(\text{nullifier}\|\text{secret})$.
  \item Send $C$ + 1 ETH to contract.  Contract checks $C$ is fresh:
\begin{lstlisting}[language=Solidity]
mapping(bytes32 => bool) public commitments;
require(!commitments[_commitment], "Already submitted");
\end{lstlisting}
  \item Contract inserts $C$ into a 20-level Merkle tree
        ($\leq 2^{20}$ deposits, no deletions):
\begin{lstlisting}[language=Solidity]
uint32 idx = _insert(_commitment);
\end{lstlisting}
  \item Contract marks $C$ seen and verifies \texttt{msg.value == 1 ETH}.
  \item Contract emits \texttt{Deposit(\_commitment, idx, block.timestamp)}.
\end{enumerate}

\section{Zero-Knowledge Proofs: zk-SNARKs}

\subsection{What Is a zk-SNARK?}

A \textbf{zk-SNARK} (Zero-Knowledge Succinct Non-Interactive Argument
of Knowledge) lets a prover convince a verifier of a statement's truth
without revealing any secret inputs.

\begin{table}[h]
\centering
\caption{zk-SNARK properties}
\begin{tabular}{ll}
\toprule
\textbf{Property} & \textbf{Meaning} \\
\midrule
Completeness    & Honest prover always generates valid proof \\
Soundness       & Infeasible to prove a false statement \\
Zero-Knowledge  & Proof reveals nothing beyond statement truth \\
Succinct        & Proof is short and fast to verify \\
Non-Interactive & No back-and-forth; single proof string \\
\bottomrule
\end{tabular}
\end{table}

\subsection{The Withdrawal Statement}

\begin{quote}
\itshape
I know a \emph{nullifier} and \emph{secret} such that:\\
(1) $\mathrm{PedersenHash}(\text{nullifier}\|\text{secret})$ is a leaf of
the Merkle tree with root \texttt{\_root}, and\\
(2) $\mathrm{PedersenHash}(\text{nullifier}) = \texttt{\_nullifierHash}$.
\end{quote}

The prover generates a proof of this statement and submits it
on-chain without revealing which commitment (leaf) it corresponds to.

\section{Withdrawal Steps}

\begin{lstlisting}[language=Solidity,caption=withdraw() signature]
function withdraw(
    bytes  calldata _proof,
    bytes32 _root,
    bytes32 _nullifierHash,
    address payable _recipient,
    address payable _relayer,
    uint256 _fee
) external;
\end{lstlisting}

\begin{enumerate}
  \item Check nullifier not yet spent (prevents double withdrawal):
\begin{lstlisting}[language=Solidity]
mapping(bytes32 => bool) public nullifierHashes;
require(!nullifierHashes[_nullifierHash], "Note already spent");
\end{lstlisting}
  \item Check \texttt{\_root} is one of the last 100 known roots (allows
        deposits while withdrawals are pending):
\begin{lstlisting}[language=Solidity]
require(isKnownRoot(_root), "Cannot find your merkle root");
\end{lstlisting}
  \item Verify SNARK proof on-chain:
\begin{lstlisting}[language=Solidity]
require(verifier.verifyProof(_proof,
    [uint256(_root), uint256(_nullifierHash), ...]),
    "Invalid withdraw proof");
\end{lstlisting}
  \item Mark nullifier spent; send ETH to recipient minus relayer fee.
\end{enumerate}

\textbf{Relayer}: Pays gas for the user (whose fresh address has no ETH)
and receives a fee deducted from the denomination.

\section{Groth16: The Proof System Used by Tornado Cash}

Tornado Cash uses the \textbf{Groth16} zk-SNARK proof system
(Jens Groth, 2016).

\subsection{Arithmetic Circuits}

Any computation can be represented as an \textbf{arithmetic circuit}
over a finite field $\mathbb{F}_p$: wires carry field elements, and
gates compute addition ($+$) and multiplication ($\times$) in
$\mathbb{F}_p$.

A \textbf{Rank-1 Constraint System (R1CS)} is a set of constraints of
the form:
\[
  (\vec{a}\cdot\vec{w})\times(\vec{b}\cdot\vec{w}) = (\vec{c}\cdot\vec{w})
\]
where $\vec{w}$ is the witness vector (all wire values: public inputs,
private inputs, and intermediate values).

\subsection{Quadratic Arithmetic Programs (QAP)}

Groth16 converts R1CS into a \textbf{QAP}: polynomials $A(x), B(x),
C(x)$ such that the constraint is satisfied iff
$A(x)\cdot B(x) - C(x) = H(x)\cdot Z(x)$ where $Z(x)=\prod_i(x-r_i)$
is the vanishing polynomial at the constraint evaluation points.

\subsection{Groth16 Proof and Verification}

\textbf{Trusted Setup}: A one-time ceremony generates
\textbf{Structured Reference Strings (SRS)} $(\sigma_1,\sigma_2)$
containing encrypted powers of a secret $\tau$.  The secret $\tau$
must be destroyed after setup (``toxic waste'').

\textbf{Proof}: Three elliptic curve points $\pi = (A, B, C) \in
G_1\times G_2\times G_1$.  Constant size ($\approx 200$ bytes)
regardless of circuit complexity.

\textbf{Verification}: Check two pairing equations:
\[
  e(A, B) = e(\alpha, \beta)\cdot e(\sum_i a_i u_i, \gamma)
           \cdot e(C, \delta)
\]
Verification takes $\approx 2$ ms on modern hardware (only 2 pairings).

\subsection{Poseidon Hash}

The original Tornado Cash uses Pedersen hash.  More modern ZK systems
use \textbf{Poseidon}, a hash function designed specifically for
arithmetic circuits:
\begin{itemize}
  \item Based on a sponge construction with a round function over
        $\mathbb{F}_p$.
  \item Far fewer constraints than SHA-256 ($\approx 240$ vs $\approx
        27\,000$ constraints per 256-bit hash).
  \item The round function: $f(x) = (x+c)^5\bmod p$ (simple
        multiplication gate).
\end{itemize}

\subsection{Trusted Setup Ceremonies}

The Groth16 setup is \textbf{circuit-specific}: a new ceremony is
needed for each circuit.  Tornado Cash used a multi-party computation
(MPC) ceremony with 1114 participants.  The system is secure as long as
\emph{at least one participant} honestly destroys their secret.

For universal setups (Plonk, STARKs), a single ceremony works for all
circuits up to a given size.

\section{The Circom Circuit}

The circuit is compiled with \textbf{circom} and a verifier contract is
generated automatically.

\begin{lstlisting}[language={},caption=withdraw.circom (simplified)]
template Withdraw(levels) {
    signal input  root;
    signal input  nullifierHash;
    signal input  recipient;  // not constrained -- prevents tampering
    signal input  relayer;
    signal input  fee;
    signal private input nullifier;
    signal private input secret;
    signal private input pathElements[levels];
    signal private input pathIndices[levels];

    // 1. Verify nullifierHash = PedersenHash(nullifier)
    component hasher = CommitmentHasher();
    hasher.nullifier <== nullifier;
    hasher.secret    <== secret;
    hasher.nullifierHash === nullifierHash;

    // 2. Verify commitment is in Merkle tree with root `root`
    component tree = MerkleTreeChecker(levels);
    tree.leaf <== hasher.commitment;
    tree.root <== root;
    for (var i = 0; i < levels; i++) {
        tree.pathElements[i] <== pathElements[i];
        tree.pathIndices[i]  <== pathIndices[i];
    }

    // Prevent optimiser from removing recipient/relayer/fee
    signal rSq; rSq <== recipient * recipient;
    signal fSq; fSq <== fee      * fee;
    signal lSq; lSq <== relayer  * relayer;
}
component main = Withdraw(20);
\end{lstlisting}

\section{OFAC Sanctions}

\begin{itemize}
  \item \textbf{Aug 8, 2022}: OFAC added Tornado Cash contracts to its
        sanctions list, alleging facilitation of ransomware money
        laundering.
  \item GitHub removed repos; three contributors suspended.
  \item Developer \textbf{Alexey Pertsev} arrested (Netherlands, Aug 2022);
        sentenced 64 months (May 2024); released Feb 2025.
  \item Developer \textbf{Roman Storm} arrested (US, Aug 2023); case ongoing.
  \item EFF and Prof. Matthew Green's advocacy led OFAC to allow
        educational use of the code.
  \item GitHub repos restored in 2023.
\end{itemize}

\begin{exambox}
  Tornado Cash illustrates the collision between financial privacy
  (a legitimate right) and AML regulations.  Core question: can
  \emph{open-source code} — not its operators — be sanctioned?
\end{exambox}


% ================================================================
\chapter{Consensus Protocols}
\label{chap:consensus}
% ================================================================

\section{State Machine Replication (SMR)}

\begin{definitionbox}
  \textbf{State Machine Replication}: $n$ nodes each maintain a local
  append-only \emph{history} (ordered list of transactions).
  \begin{itemize}
    \item \textbf{Consistency (Safety)}: All nodes agree on the same history.
    \item \textbf{Liveness}: Every valid client transaction is eventually
          added to the history.
  \end{itemize}
  If all nodes start from the same initial state and apply the same
  transactions in the same order, they reach identical final states.
\end{definitionbox}

\section{Network Models}

\begin{table}[h]
\centering
\caption{Network models for distributed consensus}
\begin{tabular}{p{3cm}p{4cm}p{5cm}}
\toprule
\textbf{Model} & \textbf{Clock} & \textbf{Message delay} \\
\midrule
Synchronous & Global, shared & Bounded: msg at step $t$ arrives by $t+1$ \\
Asynchronous & None & Unbounded (messages eventually arrive) \\
Partially synchronous & Global, shared & Unbounded before GST; $\leq\Delta$ after GST \\
\bottomrule
\end{tabular}
\end{table}

\textbf{GST} (Global Stabilisation Time): exists but is unknown in advance.

\section{Fault Types}

\begin{table}[h]
\centering
\caption{Node fault types}
\begin{tabular}{lll}
\toprule
\textbf{Fault} & \textbf{Behaviour} & \textbf{Severity} \\
\midrule
Crash     & Stops working at time $t$ & Low \\
Omission  & Drops some messages & Medium \\
Byzantine & Deviates arbitrarily; cannot forge signatures & Highest \\
\bottomrule
\end{tabular}
\end{table}

Let $f$ = Byzantine nodes, $n$ = total.

\section{SMR with Rotating Leaders (Honest Setting)}

\textbf{Assumptions}: Permissioned ($n$ known), PKI, synchronous, $f=0$.

\begin{itemize}
  \item In time step $t$: leader $= 1+(t\bmod n)$.
  \item Leader broadcasts ordered list of new transactions.
  \item Every node appends the list.
\end{itemize}

\textbf{Consistency}: All nodes receive same broadcast $\Rightarrow$
identical histories.

\textbf{Liveness}: Transaction submitted to node $i$ is included when
node $i$ becomes leader.

\section{Byzantine Broadcast (BB)}

\begin{definitionbox}
  One of $n$ nodes is the \emph{sender} (may be Byzantine) with private
  input $v^*\in V$:
  \begin{description}
    \item[Termination] Every honest node eventually halts with some $v_i\in V$.
    \item[Agreement]   All honest nodes output the same value.
    \item[Validity]    If sender is honest, all output $v^*$.
  \end{description}
\end{definitionbox}

SMR reduces to BB via rotating leaders: each time step invokes BB with
the current leader as sender.

\subsection{Protocol for $f=1$ (Cross-Checking)}

Three rounds:
\begin{enumerate}
  \item Sender broadcasts $v^*$ with signature.
  \item Each non-sender signs received message $m_i$, broadcasts $(m_i,\sigma_i)$.
        Uses $\bot$ if nothing received.
  \item Each non-sender takes majority of received messages.
\end{enumerate}

\textbf{Proposition}: Satisfies termination, agreement, validity for $f=1$.

\textbf{Fails for $f=2$}: Byzantine sender splits votes; Byzantine
non-sender amplifies the split; two groups of honest nodes disagree.

\section{The Dolev-Strong Protocol}

\begin{theorem}[Dolev-Strong, 1983]
  In the permissioned, PKI, synchronous setting, there exists a protocol
  for BB running in $f+1$ rounds that satisfies termination, agreement,
  and validity for any $f$.
\end{theorem}

\textbf{Convinced message}: node $i$ is \emph{convinced of $v$ at step
$t$} if it receives a message referencing $v$ signed by the sender and
$\geq t-1$ other distinct nodes.

\textbf{Protocol}:
\begin{itemize}
  \item Step 0: Sender sends $v^*$ + signature; outputs $v^*$.
  \item Steps $t=1,\ldots,f+1$: If convinced of $v$ and not previously,
        sign and broadcast.
  \item Final: If convinced of exactly one value $v$, output $v$; else $\bot$.
\end{itemize}

\section{Correctness of Dolev-Strong}

\begin{proof}[Proof of Agreement]
  Suppose honest nodes $i$ and $j$ both output values at the end.  We
  show they output the same value.

  For a value $v$ to be output by honest node $i$, node $i$ must have
  been convinced of $v$ at some time step $t\leq f+1$.  Being convinced
  means $i$ received a message for $v$ signed by the sender and by
  $t-1$ other distinct nodes.

  \textbf{Claim}: If honest node $i$ is first convinced of $v$ at step
  $t$, then all honest nodes are convinced of $v$ by step $t+1$.

  \textit{Reason}: At step $t$, node $i$ signs $v$ and broadcasts.
  Every honest node then receives this message signed by $i$ plus the
  $t-1$ other signatures (total: $t$ signatures including sender).
  So every honest node is convinced of $v$ by step $t+1\leq f+2$,
  before the protocol ends at step $f+1$.

  Wait — the protocol ends at step $f+1$ and conviction spreads by step
  $f+2$?  We need $t\leq f$.  If $i$ is first convinced at step $f+1$
  (the final step), it still broadcasts, but other nodes only check at
  the end of step $f+1$.  Since all messages arrive by step $f+2$ in
  the synchronous model, and all honest nodes output at the end of step
  $f+1$, we need to be more careful.

  The correct argument: honest node $i$ is convinced of $v$ at step
  $t\leq f+1$.  At step $t$, $i$ re-signs and broadcasts.  Any honest
  non-sender $j$ will receive, by step $t+1$, a message for $v$ signed
  by the sender and $t$ other nodes (including $i$).  So $j$ is
  convinced of $v$ by step $t+1\leq f+2$.  But wait, the protocol runs
  for $f+2$ steps total (steps $0, 1, \ldots, f+1$), so at step $f+1$
  (last step), $j$ receives the forwarded signatures and becomes
  convinced.  Both $i$ and $j$ are convinced of $v$, and since this
  process is symmetric, the only value any honest node can be convinced
  of is the one that spreads this way.

  If a Byzantine sender sends different values to different nodes, at
  most one chain of signatures can grow to convince honest nodes, because
  digital signatures cannot be forged.  $\square$
\end{proof}

\begin{proof}[Proof of Validity]
  If the sender is honest, it sends $v^*$ with its signature in step 0.
  Every honest non-sender receives a message signed by the sender for
  $v^*$ by step 1, so they are convinced of $v^*$.  They re-broadcast
  and all honest nodes output $v^*$.  No other value has a sender
  signature (honest sender sends only $v^*$), so no honest node is
  convinced of anything else.  $\square$
\end{proof}

\section{Impossibility Results}

\subsection{FLM (No-PKI Bound)}

\begin{theorem}[Pease--Shostak--Lamport 1980; Fischer--Lynch--Merrit 1986]
  \[
    \boxed{f\geq\frac{n}{3}
    \;\Rightarrow\;
    \text{no BB protocol without PKI}}
  \]
\end{theorem}

PKI is \emph{necessary} for Dolev-Strong.

\subsection{FLP (Asynchronous Impossibility)}

\begin{theorem}[Fischer--Lynch--Paterson, 1985]
  For $n\geq 2$, even with $f=1$, no \emph{deterministic} protocol
  solves Byzantine Agreement in the asynchronous model.
\end{theorem}

Practical blockchains escape FLP by using the \textbf{partially
synchronous} model.

\subsection{Partially Synchronous Bound}

\begin{theorem}
  There exists a deterministic BFT protocol for partially synchronous
  networks iff $f<\frac{n}{3}$.
  Safety holds at all times; liveness holds post-GST.
\end{theorem}

Tendermint and PBFT achieve this bound.

\section{Consensus Summary Table}

\begin{table}[h]
\centering
\caption{Consensus protocol comparison}
\begin{tabular}{p{2.8cm}p{2.8cm}p{2cm}p{3.5cm}p{2cm}}
\toprule
\textbf{Protocol} & \textbf{Network model} & \textbf{Max $f$} & \textbf{Guarantees} & \textbf{Used by} \\
\midrule
Rotating leaders & Sync, honest & 0 & C + L & Toy \\
Dolev-Strong & Sync, PKI & Any & T + A + V & Academic \\
PBFT & Partial sync & $<n/3$ & C + L$^*$ & HyperLedger \\
Tendermint & Partial sync & $<n/3$ & C + L$^*$ & Cosmos \\
Nakamoto & Eventual sync & $<50\%$ hash & Prob. C + L & Bitcoin \\
Casper FFG & Partial sync & $<n/3$ stake & C + L$^*$ + slash & Ethereum \\
\bottomrule
\multicolumn{5}{l}{\footnotesize $^*$Liveness only post-GST.}
\end{tabular}
\end{table}


% ================================================================
\chapter{Tendermint Protocol}
\label{chap:tendermint}
% ================================================================

\section{Overview}

\textbf{Tendermint} is a BFT consensus protocol achieving:
\begin{itemize}
  \item Consistency and eventual (post-GST) liveness.
  \item Optimal fault tolerance: $f < n/3$.
  \item Partially synchronous network, permissioned, PKI.
\end{itemize}

\textbf{Key ideas}: rotating leaders, two stages of voting per round,
restarts after timeout (if messages are delayed), weighted voting in
practice (proportional to stake).

\section{Rounds and Phases}

A \textbf{round} spans $4\Delta$ time steps:

\begin{table}[h]
\centering
\caption{Tendermint round phases}
\begin{tabular}{lll}
\toprule
\textbf{Phase} & \textbf{Start time} & \textbf{Action} \\
\midrule
1 (Propose)  & $4\Delta r$           & Leader broadcasts $(B_i, Q_i)$ \\
2 (Prevote)  & $4\Delta r+\Delta$    & Nodes cast stage-1 votes \\
3 (Precommit)& $4\Delta r+2\Delta$   & Nodes cast stage-2 votes \\
4 (Commit)   & $4\Delta r+3\Delta$   & Commit if $\geq\!\frac{2}{3}n$ stage-2 votes \\
\bottomrule
\end{tabular}
\end{table}

Each round has a unique \textbf{leader}; if no consensus, nodes move to
round $r+1$.

\section{Quorum Certificates}

Each vote has attributes: (voter $i$, block $B$, height $h$, round $r$,
stage $s$).  The triple $(h,r,s)$ is a \textbf{referendum}.

\begin{definitionbox}
  A \textbf{Quorum Certificate (QC)} for a referendum $(h,r,s)$ is a set
  of $\geq\frac{2}{3}n$ votes all for the same block in that referendum.
\end{definitionbox}

\begin{lemma}
  Any two QCs share $\geq\frac{n}{3}$ voters.
\end{lemma}

\begin{corollary}
  If $f<\frac{n}{3}$, any two QCs for the same referendum support the
  \emph{same} block (since they share at least one honest voter who
  votes at most once).
\end{corollary}

\section{Local State per Node}

Node $i$ working on block height $h_i$ maintains:
\begin{itemize}
  \item $B_i$: current belief about the correct block at height $h_i$.
  \item $Q_i$: most recent QC supporting $B_i$ (null initially).
\end{itemize}

QC ordering: $(h,r_1,s_1)>_{age}(h,r_2,s_2)$ iff $r_1>r_2$, or
$r_1=r_2$ and $s_1>s_2$.

\section{Protocol Pseudocode}

\begin{lstlisting}[language={},caption=Phase 1: Leader proposes]
if i == l then                   // i is the leader
    if l received a height-h QC newer than (B_i, Q_i) then
        B_i := B_l;  Q_i := Q_l
    end
    broadcast(B_i, Q_i)          // annotated with h_i, r, signature
end
\end{lstlisting}

\begin{lstlisting}[language={},caption=Phase 2: Stage-1 vote]
if i received (B_l, Q_l) from leader l then
    if Q_l is at least as recent as Q_i then
        B_i := B_l;  Q_i := Q_l
        broadcast(B_i, Q_i)
        broadcast stage-1-vote for B_i   // annotated h_i, r, sig
    end
end
\end{lstlisting}

\begin{lstlisting}[language={},caption=Phase 3: Stage-2 vote]
if i received >= (2n/3) round-r stage-1 votes for B then
    B_i := B;  Q_i := (those votes)      // round-r stage-1 QC
    broadcast(B_i, Q_i)
    broadcast stage-2-vote for B_i
end
\end{lstlisting}

\begin{lstlisting}[language={},caption=Phase 4: Commit]
if i received >= (2n/3) round-r stage-2 votes for B then
    B_i := B;  Q_i := (those votes)      // round-r stage-2 QC
    broadcast(B_i, Q_i)
    commit B_i to local history at height h_i
    h_i := h_i + 1
    reset B_i; reset Q_i to null
end
\end{lstlisting}

\textbf{Addendum}: If stage-2 QCs for height $h_i$ arrive before Phase~1
of round $r+1$, execute Phase~4 immediately.

\section{Proof of Consistency}

\begin{theorem}
  If $f < n/3$ and honest nodes $i,j$ both commit blocks $B,B'$ at the
  same height $h$, then $B=B'$.
\end{theorem}

\textbf{Proof sketch}:
\begin{enumerate}
  \item Let $r$ be the first round in which $>\frac{n}{3}$ honest nodes
        contribute height-$h$ stage-2 votes for a common block $B^*$.
        Call this set $S$.
  \item Any different block $B'\neq B^*$ committed in round $r+1$ would
        need a stage-2 QC in referendum $(h,r+1,2)$.
  \item Stage-2 QC requires $\geq\frac{2}{3}n$ votes; $|S|>\frac{n}{3}$
        so at least one node from $S$ must contribute.
  \item Nodes in $S$ hold $Q_i=$ stage-1 QC$(h,r,1)$ supporting $B^*$.
        They only cast stage-1 votes for a block whose QC is at least as
        recent as $Q_i$.  Stage-1 QC$(h,r+1,1)$ is more recent than
        stage-1 QC$(h,r,1)$, but it too must support $B^*$ (same argument
        applied inductively).
  \item Therefore $B^*$ is the only block that can be committed. \qed
\end{enumerate}

\section{Liveness in Tendermint (Post-GST)}

\begin{theorem}
  After GST, if $f < n/3$, Tendermint terminates (adds a block to
  history) within $O(\Delta)$ time.
\end{theorem}

\textbf{Argument}: After GST, all messages arrive within $\Delta$ time.
Consider the first round $r$ that starts after GST:
\begin{enumerate}
  \item The leader $l$ broadcasts $(B_l, Q_l)$ in Phase 1.  Since
        $f < n/3$, at least $\frac{2}{3}n$ honest nodes receive this
        message by the end of Phase 2.
  \item At least $\frac{2}{3}n$ honest nodes cast stage-1 votes for
        $B_l$ (assuming $Q_l$ is at least as recent as their local $Q$).
  \item Every honest node receives $\frac{2}{3}n$ stage-1 votes by
        Phase 3, forming a stage-1 QC, and casts a stage-2 vote.
  \item Every honest node receives $\frac{2}{3}n$ stage-2 votes by
        Phase 4 and commits.
\end{enumerate}

The ``at least as recent'' condition for stage-1 votes is subtle.  The
consistency proof shows that any block proposed by an honest leader in
round $r \geq r^*$ (where $r^*$ is the first commitment round) will be
compatible with previously committed values.

\section{Safety vs.\ Liveness in Partially Synchronous Protocols}

A fundamental tension in consensus:
\begin{itemize}
  \item \textbf{Safety} (consistency): no two honest nodes commit
        different values.  This must hold \emph{at all times}, including
        before GST.
  \item \textbf{Liveness}: transactions are eventually committed.  This
        can only be guaranteed \emph{after GST}.
\end{itemize}

\textbf{Why safety before GST?}  If the network is partitioned (two
groups cannot communicate), each group might commit a different block
(split-brain).  Tendermint avoids this by requiring $\frac{2}{3}n$
votes for commitment; in a partition where each side has $< \frac{2}{3}n$
nodes, \emph{neither side} can commit.  This sacrifices liveness
(no progress) but preserves safety (no inconsistency).

This is the essence of the \textbf{CAP theorem}: under a network
partition, a distributed system must choose between consistency and
availability.  Tendermint chooses consistency.

\section{Practical Deployment}

\begin{itemize}
  \item Cosmos SDK, Binance Smart Chain, Polygon.
  \item Nodes cast \textbf{weighted votes} (stake proportional), so
        quorums require $>\frac{2}{3}$ of total stake.
  \item See \texttt{mintscan.io} for live deployments.
  \item \textbf{Fast finality}: unlike Nakamoto consensus (probabilistic
        finality), Tendermint achieves \emph{instant} finality once a
        block is committed — it can never be reverted.
\end{itemize}


% ================================================================
\chapter{Ethereum Proof-of-Stake}
\label{chap:ethpos}
% ================================================================

\section{The Merge}

September 15, 2022: Ethereum transitioned to PoS:
\begin{itemize}
  \item Energy reduced $99.95\%$: 5.13~GW $\to$ 2.62~MW.
  \item Block time became deterministic: 12~s per slot.
  \item PoW nonce/difficulty obsoleted (set to 0).
\end{itemize}

\section{Node Architecture}

\begin{description}
  \item[Execution client] Executes transactions, maintains mempool,
       updates world state (EVM layer).
  \item[Beacon chain client] Runs the PoS algorithm; achieves consensus
       on execution payloads.
\end{description}

\section{Validators}

\begin{itemize}
  \item Must stake \textbf{32 ETH} to the deposit contract.
  \item $\approx 980\,000$ validators in March 2024.
  \item Earn attestation and proposal rewards.
  \item Slashed (penalised + forced exit) for double voting or surround
        voting.
\end{itemize}

\section{Slots and Epochs}

\begin{itemize}
  \item \textbf{Slot}: 12 seconds; one block proposed per slot.
  \item \textbf{Epoch}: 32 slots ($\approx 6.4$ minutes).
  \item Each slot: $\frac{1}{32}$ of all validators are active.
  \item One active validator per slot is the \textbf{block proposer}.
\end{itemize}

\section{Committees and Attestations}

Validators in a slot are divided into \textbf{committees}:
\begin{itemize}
  \item $128\leq$ committee size $\leq 2048$; max 64 committees.
\end{itemize}

Each validator creates an \textbf{attestation} (signed vote):

\begin{lstlisting}[language=Python,caption=Attestation data structure]
class AttestationData:
    slot:               uint64
    committee_index:    uint64
    beacon_block_root:  bytes32   # LMD GHOST vote
    source:             Checkpoint # Casper FFG source
    target:             Checkpoint # Casper FFG target

class Attestation:
    aggregation_bits:   Bitlist[MAX_VALIDATORS_PER_COMMITTEE]
    data:               AttestationData
    signature:          BLSSignature
\end{lstlisting}

\textbf{Aggregators}: subset of committee; aggregate attestations using
BLS before forwarding to the next block proposer.

\section{BLS Signatures}

\textbf{BLS} (Boneh--Lynn--Shacham) signatures support efficient
aggregation: $n$ signatures on the same message aggregate to a single
signature of the same size.

\subsection{Mathematical Foundation}

Let $G_1,G_2$ be elliptic curves of prime order $p$ with generators
$g_1,g_2$.  Let $G_T$ be a multiplicative subgroup of a finite field.

\textbf{Bilinear pairing}: $e:G_1\times G_2\to G_T$ with
\begin{enumerate}
  \item $e(g_1^\alpha,g_2^\beta)=e(g_1,g_2)^{\alpha\beta}$\quad (bilinearity)
  \item $e(g_1,g_2)\neq 1$\quad (non-degeneracy)
\end{enumerate}

\textbf{Key generation}: $sk=x$, $pk=g_1^x$.

\textbf{Sign}: $\sigma=H(m)^x\in G_2$ where $H:\{0,1\}^*\to G_2$.

\textbf{Verify}: $e(pk,H(m))\stackrel{?}{=}e(g_1,\sigma)$.

\subsection{Aggregation}

For $n$ validators with signatures $\sigma_1,\ldots,\sigma_n$ on the
same message:
\begin{align}
  \sigma_{\text{agg}} &= \prod_{i=1}^n \sigma_i \\
  pk_{\text{agg}}     &= \prod_{i=1}^n pk_i
\end{align}

Verify: $e(pk_{\text{agg}},H(m))\stackrel{?}{=}e(g_1,\sigma_{\text{agg}})$.
Constant-size proof of $n$ validators.

\section{Validator and Committee Selection}

\subsection{RANDAO: On-Chain Randomness}

Committee and block proposer assignment requires \textbf{unbiasable,
unpredictable randomness}.  Ethereum uses \textbf{RANDAO}:

\begin{enumerate}
  \item Each validator holds a private RANDAO secret $r_i$.
  \item In every block they propose, validators reveal
        $\mathrm{hash}^n(r_i)$ for decreasing $n$ (one pre-image per
        proposal).
  \item The beacon chain XORs all reveals: $\text{RANDAO} = v_1\oplus
        v_2\oplus\cdots$.
  \item At the end of each epoch, the accumulated RANDAO mix is used as
        the seed for the next epoch's committee assignments.
\end{enumerate}

\textbf{Manipulation}: A validator revealing last can choose whether
to reveal or not, gaining one bit of influence on the outcome.  This is
a known limitation; Ethereum mitigates it with VDF (Verifiable Delay
Functions) proposals for future upgrades.

\subsection{Committee Assignment Algorithm}

At the start of each epoch, validators are pseudorandomly shuffled
using \textbf{swap-or-not} (a Fisher-Yates variant efficient for
proofs):
\begin{enumerate}
  \item Take the list of active validators.
  \item For each round of the shuffle: for each index $i$, compute a
        pivot $j = \mathrm{hash}(\text{seed}\|i)$, swap $i$ with $j$
        if a coin flip (from hash) says to.
  \item The shuffled list is divided into slots and committees.
\end{enumerate}

All validators can independently compute any committee assignment given
only the RANDAO seed (no central coordinator needed).

\section{LMD GHOST Fork Choice Rule}

\textbf{LMD GHOST} = Latest Message Driven Greedy Heaviest Observed
SubTree.

\begin{itemize}
  \item \textbf{Latest}: only the most recent attestation from each
        validator counts (equivocation resistance).
  \item \textbf{Message Driven}: fork choice driven by attestations, not
        just block headers.
  \item \textbf{Greedy Heaviest Subtree}: at each fork, choose the branch
        with the most total stake weight.
\end{itemize}

\textbf{Algorithm}: starting at genesis, at each fork compute the
cumulative stake of validators whose latest attestation supports each
branch; choose the heaviest branch.

\subsection{LMD GHOST Worked Example}

\begin{example}
  Consider a block tree after a fork:
\begin{lstlisting}[language={}]
Genesis -- A -- B -- C         (branch 1)
               \-- D -- E     (branch 2)
\end{lstlisting}
  Suppose 80 validators have attested:
  \begin{itemize}
    \item 30 validators' latest attestation: block C (on branch 1)
    \item 50 validators' latest attestation: block E (on branch 2)
  \end{itemize}
  LMD GHOST: at block B's fork:
  \begin{itemize}
    \item Branch 1 (C): 30 votes
    \item Branch 2 (D/E): 50 votes
  \end{itemize}
  LMD GHOST chooses \textbf{branch 2} (heavier), so E is the head.
\end{example}

\textbf{Key insight}: LMD GHOST uses the \emph{latest} vote from each
validator.  This means a validator who voted for C and then equivocates
with a vote for E has their earlier vote \emph{replaced}, not
\emph{added}.  This prevents a simple attack where a validator
accumulates weight by voting many times.

\section{Casper FFG}

\textbf{Casper FFG} (Friendly Finality Gadget) provides economic
finality on top of LMD GHOST.

\subsection{Checkpoints}

\begin{lstlisting}[language=Python]
class Checkpoint:
    epoch: uint64    # epoch number
    root:  bytes32   # hash of first block in that epoch
\end{lstlisting}

The first slot of every epoch is a \textbf{checkpoint}.

\subsection{FFG Votes}

An FFG vote is a link $s\to t$ where:
\begin{itemize}
  \item $s$ = highest justified checkpoint the validator knows (\emph{source})
  \item $t$ = current epoch checkpoint (\emph{target})
  \item $s$ must be an ancestor of $t$
\end{itemize}

A \textbf{supermajority link} $s\to t$: $\geq\frac{2}{3}$ of validators
voted for it.

\subsection{Justification}

If $c_1$ is justified and $c_1\to c_2$ is a supermajority link,
$c_2$ becomes \textbf{justified}.

\subsection{Finalization}

If justified $c_1\to c_2$ is a supermajority link and $c_2$ is the
\emph{direct child epoch} of $c_1$, then $c_1$ is \textbf{finalized}.

\begin{keyinsight}
  Once a checkpoint is finalized, it is irreversible as long as
  $<\frac{1}{3}$ of the stake is Byzantine.
  If two conflicting checkpoints are finalized, at least $\frac{1}{3}$
  of the stake \emph{can be identified} and slashed
  (\textbf{accountable safety}).
\end{keyinsight}

\section{Casper FFG Slashing Rules}

\subsection{Rule 1 — No Double Vote}

A validator must not cast two distinct votes $s_1\to t_1$ and
$s_2\to t_2$ with $h(t_1)=h(t_2)$ (same target epoch height).

\subsection{Rule 2 — No Surround Vote}

A validator must not cast $s_1\to t_1$ and $s_2\to t_2$ such that:
\[
  h(s_1) < h(s_2) < h(t_2) < h(t_1)
\]
(one vote ``surrounds'' the other).

\section{Accountable Safety}

\begin{theorem}
  If $<\frac{1}{3}$ of validators are Byzantine and two conflicting
  checkpoints $a_m,b_n$ ($m<n$) are both finalized, then at least
  $\frac{1}{3}$ of total stake can be proven to have violated Rule~1
  or Rule~2, and hence slashed.
\end{theorem}

\textbf{Proof sketch}: There is a chain of supermajority links from
genesis to $b_n$.  The finalization of $a_m$ requires a supermajority
link $a_m\to a_{m+1}$.  The pair $(a_m,a_{m+1})$ must fall between two
consecutive links $b_{j-1}\to b_j$ in the chain to $b_n$.  Any
validator who voted for both $a_m\to a_{m+1}$ and a link in the $b_n$
chain violated Rule~1 (same target epoch) or Rule~2 (surround vote).
Since two QCs overlap by $>\frac{1}{3}$ stake, enough violators are
identifiable.  $\square$

\section{Ethereum PoS Parameter Summary}

\begin{table}[h]
\centering
\caption{Ethereum PoS key parameters}
\begin{tabular}{ll}
\toprule
\textbf{Parameter} & \textbf{Value} \\
\midrule
Slot duration          & 12 seconds \\
Epoch                  & 32 slots $\approx 6.4$ min \\
Validator deposit      & 32 ETH \\
Validators per slot    & $\approx n/32$ \\
Committee size         & 128--2048 \\
Fault tolerance        & $f<\frac{1}{3}$ of total stake \\
Finality latency       & $\approx$ 2 epochs ($\approx 12.8$ min) \\
Fork choice            & LMD GHOST \\
Finality gadget        & Casper FFG \\
Signature scheme       & BLS12-381 \\
\bottomrule
\end{tabular}
\end{table}

\begin{exambox}
  \textbf{Ethereum PoS vs.\ Tendermint}: Both achieve $f<n/3$ BFT in
  the partially synchronous model.  Tendermint uses one-vote-per-node;
  Ethereum uses stake-weighted BLS aggregation, enabling participation
  of $\sim\!10^6$ validators at manageable bandwidth.  LMD GHOST handles
  fork choice; Casper FFG adds economic finality with slashing.
\end{exambox}


% ================================================================
\appendix
\chapter{Mathematical Quick Reference}
% ================================================================

\section{Group Theory}

\begin{table}[h]
\centering
\caption{Group theory identities}
\begin{tabular}{ll}
\toprule
\textbf{Statement} & \textbf{Formula} \\
\midrule
Fermat's Little Thm & $a^{p-1}\equiv 1\pmod p$ ($p$ prime, $p\nmid a$) \\
Modular inverse     & $a^{-1}\equiv a^{p-2}\pmod p$ \\
Euler's Criterion   & $a^{(p-1)/2}\equiv\pm 1\pmod p$ \\
Lagrange's Theorem  & $|H|\ \mid\ |G|$ for $H\leq G$ \\
\bottomrule
\end{tabular}
\end{table}

\section{Elliptic Curve Operations}

\begin{table}[h]
\centering
\caption{ECC quick reference}
\begin{tabular}{p{3.5cm}p{6cm}p{2cm}}
\toprule
\textbf{Operation} & \textbf{Formula} & \textbf{Complexity} \\
\midrule
Addition $P+Q$ & $\lambda=\frac{y_2-y_1}{x_2-x_1}$; $x_3=\lambda^2-x_1-x_2$; $y_3=\lambda(x_1-x_3)-y_1$ & $O(1)$ \\
Doubling $2P$ & $\lambda=\frac{3x_1^2+a}{2y_1}$; $x_3=\lambda^2-2x_1$; $y_3=\lambda(x_1-x_3)-y_1$ & $O(1)$ \\
Scalar mult $kP$ & double-and-add & $O(\log k)$ \\
ECDLP & find $k$ given $G,kG$ & $O(\sqrt{n})$ \\
\bottomrule
\end{tabular}
\end{table}

\section{Hash Function Summary}

\begin{table}[h]
\centering
\caption{Hash functions used in Bitcoin and Ethereum}
\begin{tabular}{llll}
\toprule
\textbf{Function} & \textbf{Output} & \textbf{Preimage} & \textbf{Collision} \\
\midrule
SHA-256       & 256 bit & $2^{256}$ & $2^{128}$ \\
RIPEMD-160    & 160 bit & $2^{160}$ & $2^{80}$  \\
SHA256d       & 256 bit & $2^{256}$ & $2^{128}$ \\
Keccak-256    & 256 bit & $2^{256}$ & $2^{128}$ \\
\bottomrule
\end{tabular}
\end{table}


% ================================================================
\chapter{Glossary}
% ================================================================

\begin{description}[leftmargin=2.8cm,style=sameline]
  \item[ASIC]       Application-Specific Integrated Circuit; custom mining chip.
  \item[BFT]        Byzantine Fault Tolerant: works with $<\frac{1}{3}$ faulty nodes.
  \item[BLS]        Boneh--Lynn--Shacham pairing-based signature scheme.
  \item[Calldata]   Read-only input data of a transaction or call.
  \item[CEI]        Checks-Effects-Interactions: secure Solidity pattern.
  \item[ECDLP]      Elliptic Curve Discrete Logarithm Problem.
  \item[ECDSA]      Elliptic Curve Digital Signature Algorithm.
  \item[EOA]        Externally Owned Account; controlled by a private key.
  \item[EVM]        Ethereum Virtual Machine.
  \item[Gas]        Unit of EVM computation.
  \item[GST]        Global Stabilisation Time.
  \item[LMD GHOST]  Latest Message Driven Greedy Heaviest Observed SubTree.
  \item[Merkle tree] Binary hash tree enabling $O(\log n)$ membership proofs.
  \item[MPT]        Merkle Patricia Trie; Ethereum's state storage.
  \item[NFT]        Non-Fungible Token (ERC-721).
  \item[PoS]        Proof of Stake.
  \item[PoW]        Proof of Work.
  \item[QC]         Quorum Certificate: $\geq\frac{2}{3}n$ votes for same block.
  \item[RLP]        Recursive Length Prefix encoding.
  \item[SNARK]      Succinct Non-Interactive Argument of Knowledge.
  \item[UTXO]       Unspent Transaction Output.
  \item[Wei]        Smallest ETH unit; $10^{18}\ \text{wei}=1\ \text{ETH}$.
  \item[zk-SNARK]   Zero-Knowledge SNARK.
\end{description}

\end{document}
